\begin{thebibliography}{999}
\addcontentsline{toc}{chapter}{Bibliography}
\small
\bibitem{abatewhitt2006} J. Abate and W. Whitt. A unified framework for numerically inverting Laplace transforms. INFORMS J. Comput. 18 (2006).
\bibitem{abb2017} L. P. Arguin, D. Belius and P. Bourgade. Maximum of the characteristic polynomial of random unitary matrices. Comm. Math. Phys. 349 (2017).
\bibitem{abbrs2019} L. P. Arguin, D. Belius, P. Bourgade, M. Radziwi\l\l\ and K. Soundararajan. Maximum of the Riemann zeta function on a short interval of the critical line. Comm. Pure Appl. Math. 72 (2019).
\bibitem{abh2017} L. P. Arguin, D. Belius and A. J. Harper. Maxima of a randomized Riemann zeta function, and branching random walks. Ann. Appl. Probab. 27 (2017).
\bibitem{abl2021} P. Autissier, D. Bonolis and Y. Lamzouri. The distribution of the maximum of partial sums of Kloosterman sums and other trace functions. Compos. Math. 157 (2021).
\bibitem{abr2020} L. P. Arguin, P. Bourgade and M. Radziwi\l\l. The Fyodorov Hiary Keating conjecture. I. Preprint, arXiv 2007.00988 (2020).
\bibitem{abr2023} L. P. Arguin, P. Bourgade and M. Radziwi\l\l. The Fyodorov Hiary Keating conjecture. II. Preprint, arXiv 2307.00982 (2023).
\bibitem{accghlnstt2023} P. B. Allen, F. Calegari, A. Caraiani, T. Gee, D. Helm, B. V. Le Hung, J. Newton, P. Scholze, R. Taylor and J. A. Thorne. Potential automorphy over CM fields. Ann. of Math. 197 (2023).
\bibitem{addarioberryreed2009} L. Addario Berry and B. Reed. Minima in branching random walks. Ann. Probab. 37 (2009).
\bibitem{aidekon2013} E. A\"id\'ekon. Convergence in law of the minimum of a branching random walk. Ann. Probab. 41 (2013).
\bibitem{akhiezer1965} N. I. Akhiezer. The Classical Moment Problem and Some Related Questions in Analysis. Oliver and Boyd (1965).
\bibitem{akiyamatanigawa1999} S. Akiyama and Y. Tanigawa. Calculation of values of L functions associated to elliptic curves. Math. Comp. 68 (1999).
\bibitem{aldous1990} D. J. Aldous. The random walk construction of uniform spanning trees and uniform labelled trees. SIAM J. Discrete Math. 3 (1990).
\bibitem{aldousbrown1992} D. J. Aldous and M. Brown. Inequalities for rare events in time reversible Markov chains I. In Stochastic Inequalities, IMS Lecture Notes Monograph Series 22 (1992).
\bibitem{aldousfill} D. Aldous and J. A. Fill. Reversible Markov Chains and Random Walks on Graphs. Unfinished monograph, available online (2002).
\bibitem{alladi1977} K. Alladi. Duality between prime factors and an application to the prime number theorem for arithmetic progressions. J. Number Theory 9 (1977).
\bibitem{allen2023} P. B. Allen, F. Calegari, A. Caraiani, T. Gee, D. Helm, B. V. Le Hung, J. Newton, P. Scholze, R. Taylor and J. A. Thorne. Potential automorphy over CM fields. Ann. of Math. 197 (2023).
\bibitem{alon2021} N. Alon. Explicit expanders of every degree and size. Combinatorica 41 (2021).
\bibitem{amdeberhan2012} T. Amdeberhan, O. Espinosa, I. Gonzalez, M. Harrison, V. H. Moll and A. Straub. Ramanujan's Master Theorem. Ramanujan J. 29 (2012).
\bibitem{antoniak1974} C. E. Antoniak. Mixtures of Dirichlet processes with applications to Bayesian nonparametric problems. Ann. Statist. 2 (1974).
\bibitem{arangopinerokeliherkeyes2022} S. Arango Pi\~neros, D. Keliher and C. Keyes. Mertens' theorem for Chebotarev sets. Int. J. Number Theory 18 (2022).
\bibitem{armitage1972} J. V. Armitage. Zeta functions with a zero at $s=\frac12$. Invent. Math. 15 (1972).
\bibitem{arratiabarbourtavare2003} R. Arratia, A. D. Barbour and S. Tavar\'e. Logarithmic Combinatorial Structures: a Probabilistic Approach. European Mathematical Society (2003).
\bibitem{arratiakochmanmiller2014} R. Arratia, F. Kochman and V. S. Miller. Extensions of Billingsley's theorem via multi intensities. Preprint, arXiv 1401.1555 (2014).
\bibitem{arthur2003} J. Arthur. The principle of functoriality. Bull. Amer. Math. Soc. 40 (2003).
\bibitem{arthur2013} J. Arthur. The Endoscopic Classification of Representations: Orthogonal and Symplectic Groups. AMS Colloquium Publications 61 (2013).
\bibitem{arthurclozel1989} J. Arthur and L. Clozel. Simple Algebras, Base Change, and the Advanced Theory of the Trace Formula. Annals of Mathematics Studies 120, Princeton University Press (1989).
\bibitem{artin1924} E. Artin. \"Uber eine neue Art von L Reihen. Abh. Math. Sem. Univ. Hamburg 3 (1924).
\bibitem{artin1927} E. Artin. Beweis des allgemeinen Reziprozit\"atsgesetzes. Abh. Math. Sem. Hamburg 5 (1927).
\bibitem{atkinlehner1970} A. O. L. Atkin and J. Lehner. Hecke operators on $\Gamma_0(m)$. Math. Ann. 185 (1970).
\bibitem{backlund1914} R. J. Backlund. Sur les z\'eros de la fonction $\zeta(s)$ de Riemann. C. R. Acad. Sci. Paris 158 (1914).
\bibitem{bagchi1981} B. Bagchi. The statistical behaviour and universality properties of the Riemann zeta function and other allied Dirichlet series. PhD thesis, Indian Statistical Institute, Calcutta (1981).
\bibitem{baileykeating2019} E. C. Bailey and J. P. Keating. On the moments of the moments of the characteristic polynomials of random unitary matrices. Comm. Math. Phys. 371 (2019).
\bibitem{bakernorine2007} M. Baker and S. Norine. Riemann Roch and Abel Jacobi theory on a finite graph. Adv. Math. 215 (2007).
\bibitem{banaszakkedlaya2015} G. Banaszak and K. S. Kedlaya. An algebraic Sato Tate group and Sato Tate conjecture. Indiana Univ. Math. J. 64 (2015).
\bibitem{bardestanifreiberg2013} M. Bardestani and T. Freiberg. Mertens's theorem for splitting primes and more. Preprint, arXiv 1309.7482 (2013).
\bibitem{barndorffnielsen1980} O. E. Barndorff Nielsen. Conditionality resolutions. Biometrika 67 (1980).
\bibitem{barndorffnielsen1983} O. E. Barndorff Nielsen. On a formula for the distribution of the maximum likelihood estimator. Biometrika 70 (1983).
\bibitem{barndorffnielsencox1994} O. E. Barndorff Nielsen and D. R. Cox. Inference and Asymptotics. Chapman and Hall (1994).
\bibitem{bass1992} H. Bass. The Ihara Selberg zeta function of a tree lattice. Internat. J. Math. 3 (1992).
\bibitem{bblm2006} E. Bogomolny, O. Bohigas, P. Leboeuf and A. G. Monastra. On the spacing distribution of the Riemann zeros: corrections to the asymptotic result. J. Phys. A 39 (2006).
\bibitem{bcdt2001} C. Breuil, B. Conrad, F. Diamond and R. Taylor. On the modularity of elliptic curves over $\mathbb{Q}$: wild 3 adic exercises. J. Amer. Math. Soc. 14 (2001).
\bibitem{bcgp2021} G. Boxer, F. Calegari, T. Gee and V. Pilloni. Abelian surfaces over totally real fields are potentially modular. Publ. Math. Inst. Hautes \'Etudes Sci. 134 (2021).
\bibitem{bellman1961} R. Bellman. A Brief Introduction to Theta Functions. Holt, Rinehart and Winston (1961).
\bibitem{benjaminischramm2001} I. Benjamini and O. Schramm. Recurrence of distributional limits of finite planar graphs. Electron. J. Probab. 6 (2001).
\bibitem{berard1980} P. H. B\'erard. Spectres et groupes cristallographiques I: domaines euclidiens. Invent. Math. 58 (1980).
\bibitem{berger1985} J. O. Berger. Statistical Decision Theory and Bayesian Analysis. Second edition, Springer (1985).
\bibitem{bergerbernardo1992} J. O. Berger and J. M. Bernardo. On the development of reference priors. In Bayesian Statistics 4, Oxford University Press (1992).
\bibitem{bergeronvenkatesh2013} N. Bergeron and A. Venkatesh. The asymptotic growth of torsion homology for arithmetic groups. J. Inst. Math. Jussieu 12 (2013).
\bibitem{berk1966} R. H. Berk. Limiting behavior of posterior distributions when the model is incorrect. Ann. Math. Statist. 37 (1966).
\bibitem{bernardo1979} J. M. Bernardo. Reference posterior distributions for Bayesian inference. J. R. Statist. Soc. B 41 (1979).
\bibitem{bernstein1929} S. Bernstein. Sur les fonctions absolument monotones. Acta Math. 52 (1929).
\bibitem{bertoin1996} J. Bertoin. L\'evy Processes. Cambridge Tracts in Mathematics 121, Cambridge University Press (1996).
\bibitem{besicovitch1932} A. S. Besicovitch. Almost Periodic Functions. Cambridge University Press (1932).
\bibitem{bew1998} B. C. Berndt, R. J. Evans and K. S. Williams. Gauss and Jacobi Sums. Wiley (1998).
\bibitem{bhargava2005} M. Bhargava. The density of discriminants of quartic rings and fields. Ann. of Math. 162 (2005).
\bibitem{bhargavashankar2013} M. Bhargava and A. Shankar. The average size of the 5 Selmer group of elliptic curves is 6, and the average rank is less than 1. Preprint, arXiv 1312.7859 (2013).
\bibitem{bhargavashankar2015} M. Bhargava and A. Shankar. Binary quartic forms having bounded invariants, and the boundedness of the average rank of elliptic curves. Ann. of Math. 181 (2015).
\bibitem{bhargavashankar2015b} M. Bhargava and A. Shankar. Ternary cubic forms having bounded invariants, and the existence of a positive proportion of elliptic curves having rank 0. Ann. of Math. 181 (2015).
\bibitem{bhkssw2016} J. S. Balakrishnan, W. Ho, N. Kaplan, S. Spicer, W. Stein and J. Weigandt. Databases of elliptic curves ordered by height and distributions of Selmer groups and ranks. LMS J. Comput. Math. 19 (2016).
\bibitem{biane2009} P. Biane. Matrix valued Brownian motion and a paper by P\'olya. In S\'eminaire de Probabilit\'es XLII, Lecture Notes in Math. 1979, Springer (2009).
\bibitem{biggs1997} N. Biggs. Algebraic potential theory on graphs. Bull. London Math. Soc. 29 (1997).
\bibitem{billingsley1968} P. Billingsley. Convergence of Probability Measures. Wiley (1968).
\bibitem{billingsley1972} P. Billingsley. On the distribution of large prime divisors. Period. Math. Hungar. 2 (1972).
\bibitem{billingsley1995} P. Billingsley. Probability and Measure. Third edition, Wiley (1995).
\bibitem{binghamgoldieteugels1987} N. H. Bingham, C. M. Goldie and J. L. Teugels. Regular Variation. Cambridge University Press (1987).
\bibitem{birch1968} B. J. Birch. How the number of points of an elliptic curve over a fixed prime field varies. J. London Math. Soc. 43 (1968).
\bibitem{bklpr2015} M. Bhargava, D. M. Kane, H. W. Lenstra, B. Poonen and E. Rains. Modeling the distribution of ranks, Selmer groups, and Shafarevich Tate groups of elliptic curves. Camb. J. Math. 3 (2015).
\bibitem{blackwellmacqueen1973} D. Blackwell and J. B. MacQueen. Ferguson distributions via P\'olya urn schemes. Ann. Statist. 1 (1973).
\bibitem{blght2011} T. Barnet Lamb, D. Geraghty, M. Harris and R. Taylor. A family of Calabi Yau varieties and potential automorphy II. Publ. Res. Inst. Math. Sci. 47 (2011).
\bibitem{bmsw2007} B. Bektemirov, B. Mazur, W. Stein and M. Watkins. Average ranks of elliptic curves: tension between data and conjecture. Bull. Amer. Math. Soc. 44 (2007).
\bibitem{bogomolnykeating1996} E. B. Bogomolny and J. P. Keating. Gutzwiller's trace formula and spectral statistics: beyond the diagonal approximation. Phys. Rev. Lett. 77 (1996).
\bibitem{bohrjessen1930} H. Bohr and B. Jessen. \"Uber die Werteverteilung der Riemannschen Zetafunktion I, II. Acta Math. 54 (1930), 58 (1932).
\bibitem{bombieri1974} E. Bombieri. Counting points on curves over finite fields (d'apr\`es S. A. Stepanov). S\'eminaire Bourbaki 430, Lecture Notes in Math. 383 (1974).
\bibitem{bombierihejhal1995} E. Bombieri and D. A. Hejhal. On the distribution of zeros of linear combinations of Euler products. Duke Math. J. 80 (1995).
\bibitem{bondesson1992} L. Bondesson. Generalized Gamma Convolutions and Related Classes of Distributions and Densities. Lecture Notes in Statistics 76, Springer (1992).
\bibitem{bondesson2015} L. Bondesson. A class of probability distributions that is closed with respect to addition as well as multiplication of independent random variables. J. Theoret. Probab. 28 (2015).
\bibitem{booker2006} A. R. Booker. Artin's conjecture, Turing's method, and the Riemann hypothesis. Experiment. Math. 15 (2006).
\bibitem{bookerstrombergssonvenkatesh2006} A. R. Booker, A. Str\"ombergsson and A. Venkatesh. Effective computation of Maass cusp forms. Int. Math. Res. Not. (2006).
\bibitem{bordenave2020} C. Bordenave. A new proof of Friedman's second eigenvalue theorem and its extension to random lifts. Ann. Sci. \'Ec. Norm. Sup\'er. 53 (2020).
\bibitem{borel1979} A. Borel. Automorphic L functions. Proc. Sympos. Pure Math. 33, Part 2, Amer. Math. Soc. (1979).
\bibitem{boschsimon2016} P. Bosch and T. Simon. A proof of Bondesson's conjecture on stable densities. Ark. Mat. 54 (2016).
\bibitem{bourgade2010} P. Bourgade. Mesoscopic fluctuations of the zeta zeros. Probab. Theory Related Fields 148 (2010).
\bibitem{bovier2017} A. Bovier. Gaussian Processes on Trees: From Spin Glasses to Branching Brownian Motion. Cambridge Studies in Advanced Mathematics 163, Cambridge University Press (2017).
\bibitem{bpy2001} P. Biane, J. Pitman and M. Yor. Probability laws related to the Jacobi theta and Riemann zeta functions, and Brownian excursions. Bull. Amer. Math. Soc. 38 (2001).
\bibitem{bramson1978} M. D. Bramson. Maximal displacement of branching Brownian motion. Comm. Pure Appl. Math. 31 (1978).
\bibitem{brauer1947} R. Brauer. On Artin's L series with general group characters. Ann. of Math. 48 (1947).
\bibitem{broder1989} A. Broder. Generating random spanning trees. In Proc. 30th IEEE Symp. Foundations of Computer Science (1989).
\bibitem{broughan2009} K. A. Broughan. Extension of the Riemann $\xi$ function's logarithmic derivative positivity region to near the critical strip. Canad. Math. Bull. 52 (2009).
\bibitem{bruggeman1978} R. W. Bruggeman. Fourier coefficients of cusp forms. Invent. Math. 45 (1978).
\bibitem{brumer1992} A. Brumer. The average rank of elliptic curves I. Invent. Math. 109 (1992).
\bibitem{bsd1965} B. J. Birch and H. P. F. Swinnerton Dyer. Notes on elliptic curves II. J. Reine Angew. Math. 218 (1965).
\bibitem{bst2013} M. Bhargava, A. Shankar and J. Tsimerman. On the Davenport Heilbronn theorems and second order terms. Invent. Math. 193 (2013).
\bibitem{bsz2014} M. Bhargava, C. Skinner and W. Zhang. A majority of elliptic curves over $\mathbb{Q}$ satisfy the Birch and Swinnerton Dyer conjecture. Preprint, arXiv 1407.1826 (2014).
\bibitem{bump1997} D. Bump. Automorphic Forms and Representations. Cambridge University Press (1997).
\bibitem{burgess1963} D. A. Burgess. On character sums and L series II. Proc. London Math. Soc. 13 (1963).
\bibitem{calegari2011} F. Calegari. Even Galois representations and the Fontaine Mazur conjecture. Invent. Math. 185 (2011).
\bibitem{calegari2012} F. Calegari. Even Galois representations and the Fontaine Mazur conjecture II. J. Amer. Math. Soc. 25 (2012).
\bibitem{calegarigeraghty2018} F. Calegari and D. Geraghty. Modularity lifting beyond the Taylor Wiles method. Invent. Math. 211 (2018); correction, Invent. Math. 227 (2022).
\bibitem{calegarivenkatesh2019} F. Calegari and A. Venkatesh. A Torsion Jacquet Langlands Correspondence. Ast\'erisque 409 (2019).
\bibitem{cardy1992} J. L. Cardy. Critical percolation in finite geometries. J. Phys. A 25 (1992).
\bibitem{cardy2002} J. L. Cardy. Crossing formulae for critical percolation in an annulus. J. Phys. A 35 (2002).
\bibitem{cardy2006} J. L. Cardy. The O(n) model on the annulus. J. Stat. Phys. 125 (2006).
\bibitem{castryckdecru2023} W. Castryck and T. Decru. An efficient key recovery attack on SIDH. In Advances in Cryptology, Eurocrypt 2023, Lecture Notes in Comput. Sci. 14008 (2023).
\bibitem{cdf1997} J. B. Conrey, W. Duke and D. W. Farmer. The distribution of the eigenvalues of Hecke operators. Acta Arith. 78 (1997).
\bibitem{cfkrs2005} J. B. Conrey, D. Farmer, J. Keating, M. Rubinstein and N. Snaith. Integral moments of L functions. Proc. London Math. Soc. 91 (2005).
\bibitem{cfz2008} J. B. Conrey, D. Farmer and M. Zirnbauer. Autocorrelation of ratios of L functions. Commun. Number Theory Phys. 2 (2008).
\bibitem{cgl2009} D. X. Charles, E. Z. Goren and K. E. Lauter. Cryptographic hash functions from expander graphs. J. Cryptology 22 (2009).
\bibitem{chandra1996} A. K. Chandra, P. Raghavan, W. L. Ruzzo, R. Smolensky and P. Tiwari. The electrical resistance of a graph captures its commute and cover times. Comput. Complexity 6 (1996).
\bibitem{changeaves1990} T. Chang and D. Eaves. Reference priors for the orbit in a group model. Ann. Statist. 18 (1990).
\bibitem{chavel1984} I. Chavel. Eigenvalues in Riemannian Geometry. Academic Press (1984).
\bibitem{chebotarev1926} N. Tschebotareff. Die Bestimmung der Dichtigkeit einer Menge von Primzahlen, welche zu einer gegebenen Substitutionsklasse geh\"oren. Math. Ann. 95 (1926).
\bibitem{cht2008} L. Clozel, M. Harris and R. Taylor. Automorphy for some $\ell$ adic lifts of automorphic mod $\ell$ Galois representations. Publ. Math. IHES 108 (2008).
\bibitem{ciesielskitaylor1962} Z. Ciesielski and S. J. Taylor. First passage times and sojourn times for Brownian motion in space and the exact Hausdorff measure of the sample path. Trans. Amer. Math. Soc. 103 (1962).
\bibitem{cifarelliregazzini1990} D. M. Cifarelli and E. Regazzini. Distribution functions of means of a Dirichlet process. Ann. Statist. 18 (1990).
\bibitem{cmn2018} R. Chhaibi, T. Madaule and J. Najnudel. On the maximum of the C$\beta$E field. Duke Math. J. 167 (2018).
\bibitem{cmsv2019} E. Costa, N. Mascot, J. Sijsling and J. Voight. Rigorous computation of the endomorphism ring of a Jacobian. Math. Comp. 88 (2019).
\bibitem{cogdellps1994} J. W. Cogdell and I. I. Piatetski Shapiro. Converse theorems for $\mathrm{GL}_n$. Publ. Math. IHES 79 (1994).
\bibitem{cohen2016} M. B. Cohen. Ramanujan graphs in polynomial time. Proc. IEEE Symposium on Foundations of Computer Science (2016).
\bibitem{cohenlenstra1984} H. Cohen and H. W. Lenstra. Heuristics on class groups of number fields. In Number Theory, Noordwijkerhout 1983, Lecture Notes in Math. 1068, Springer (1984).
\bibitem{cohenmartinet1990} H. Cohen and J. Martinet. \'Etude heuristique des groupes de classes des corps de nombres. J. Reine Angew. Math. 404 (1990).
\bibitem{colemanedixhoven1998} R. F. Coleman and B. Edixhoven. On the semi simplicity of the $U_p$ operator on modular forms. Math. Ann. 310 (1998).
\bibitem{conrad2005} K. Conrad. Partial Euler products on the critical line. Canad. J. Math. 57 (2005).
\bibitem{conrey1989} J. B. Conrey. More than two fifths of the zeros of the Riemann zeta function are on the critical line. J. Reine Angew. Math. 399 (1989).
\bibitem{conreyfarmer1995} J. B. Conrey and D. W. Farmer. An extension of Hecke's converse theorem. Internat. Math. Res. Notices (1995).
\bibitem{conreyghosh1998} J. B. Conrey and A. Ghosh. A conjecture for the sixth power moment of the Riemann zeta function. Internat. Math. Res. Notices (1998).
\bibitem{conreygonek2001} J. B. Conrey and S. M. Gonek. High moments of the Riemann zeta function. Duke Math. J. 107 (2001).
\bibitem{conreysnaith2007} J. B. Conrey and N. C. Snaith. Applications of the $L$ functions ratios conjectures. Proc. London Math. Soc. 94 (2007).
\bibitem{cooperfrieze2005} C. Cooper and A. Frieze. The cover time of random regular graphs. SIAM J. Discrete Math. 18 (2005).
\bibitem{csordassmithvarga1994} G. Csordas, W. Smith and R. S. Varga. Lehmer pairs of zeros, the de Bruijn Newman constant $\Lambda$, and the Riemann hypothesis. Constr. Approx. 10 (1994).
\bibitem{curryschoenberg1966} H. B. Curry and I. J. Schoenberg. On P\'olya frequency functions IV: the fundamental spline functions and their limits. J. Analyse Math. 17 (1966).
\bibitem{daniels1954} H. E. Daniels. Saddlepoint approximations in statistics. Ann. Math. Statist. 25 (1954).
\bibitem{dattaghosh1995} G. S. Datta and J. K. Ghosh. Noninformative priors for maximal invariant parameter in group models. Test 4 (1995).
\bibitem{dattaghosh1996} G. S. Datta and M. Ghosh. On the invariance of noninformative priors. Ann. Statist. 24 (1996).
\bibitem{davenport2000} H. Davenport. Multiplicative Number Theory, third edition, revised by H. L. Montgomery. Graduate Texts in Mathematics 74, Springer (2000).
\bibitem{davenportheilbronn1936} H. Davenport and H. Heilbronn. On the zeros of certain Dirichlet series. J. London Math. Soc. 11 (1936).
\bibitem{davenportheilbronn1971} H. Davenport and H. Heilbronn. On the density of discriminants of cubic fields II. Proc. Roy. Soc. London Ser. A 322 (1971).
\bibitem{davidpappalardi1999} C. David and F. Pappalardi. Average Frobenius distributions of elliptic curves. Internat. Math. Res. Notices (1999).
\bibitem{dawidstonezidek1973} A. P. Dawid, M. Stone and J. V. Zidek. Marginalization paradoxes in Bayesian and structural inference. J. R. Statist. Soc. B 35 (1973).
\bibitem{dawsey2017} M. L. Dawsey. A new formula for Chebotarev densities. Res. Number Theory 3 (2017).
\bibitem{ddt1995} H. Darmon, F. Diamond and R. Taylor. Fermat's Last Theorem. In Current Developments in Mathematics, International Press (1995).
\bibitem{deboor2005} C. de Boor. Divided differences. Surveys in Approximation Theory 1 (2005).
\bibitem{debruijn1950} N. G. de Bruijn. The roots of trigonometric integrals. Duke Math. J. 17 (1950).
\bibitem{definetti1937} B. de Finetti. La pr\'evision: ses lois logiques, ses sources subjectives. Ann. Inst. H. Poincar\'e 7 (1937).
\bibitem{delange1971} H. Delange. Sur des formules de Atle Selberg. Acta Arith. 19 (1971).
\bibitem{delaunay2001} C. Delaunay. Heuristics on Tate Shafarevitch groups of elliptic curves defined over $\mathbb{Q}$. Experiment. Math. 10 (2001).
\bibitem{delavalleepoussin1896} C. J. de la Vall\'ee Poussin. Recherches analytiques sur la th\'eorie des nombres premiers. Ann. Soc. Sci. Bruxelles 20 (1896).
\bibitem{delfsgalbraith2016} C. Delfs and S. D. Galbraith. Computing isogenies between supersingular elliptic curves over $\mathbb{F}_p$. Des. Codes Cryptogr. 78 (2016).
\bibitem{deligne1971} P. Deligne. Formes modulaires et repr\'esentations $\ell$ adiques. S\'eminaire Bourbaki 355, Lecture Notes in Math. 179 (1971).
\bibitem{deligne1974} P. Deligne. La conjecture de Weil I. Publ. Math. IHES 43 (1974).
\bibitem{deligne1980} P. Deligne. La conjecture de Weil II. Publ. Math. IHES 52 (1980).
\bibitem{deligneserre1974} P. Deligne and J. P. Serre. Formes modulaires de poids 1. Ann. Sci. \'Ecole Norm. Sup. 7 (1974).
\bibitem{dembozeitouni1998} A. Dembo and O. Zeitouni. Large Deviations Techniques and Applications. Second edition, Springer (1998).
\bibitem{demontrichertapiathompson1975} G. F. de Montricher, R. A. Tapia and J. R. Thompson. Nonparametric maximum likelihood estimation of probability densities by penalty function methods. Ann. Statist. 3 (1975).
\bibitem{derickx2013} M. Derickx, M. van Hoeij and J. Zeng. Computing Galois representations and equations for modular curves $X_H(\ell)$. Preprint, arXiv 1312.6819 (2013).
\bibitem{deshouillersiwaniec1982} J. M. Deshouillers and H. Iwaniec. Kloosterman sums and Fourier coefficients of cusp forms. Invent. Math. 70 (1982).
\bibitem{dettestudden1997} H. Dette and W. J. Studden. The Theory of Canonical Moments with Applications in Statistics, Probability, and Analysis. Wiley (1997).
\bibitem{deuring1941} M. Deuring. Die Typen der Multiplikatorenringe elliptischer Funktionenk\"orper. Abh. Math. Sem. Hansischen Univ. 14 (1941).
\bibitem{deuring1953} M. Deuring. Die Zetafunktion einer algebraischen Kurve vom Geschlechte Eins. Nachr. Akad. Wiss. G\"ottingen (1953).
\bibitem{dhkms2012} E. Due\~nez, D. K. Huynh, J. P. Keating, S. J. Miller and N. C. Snaith. A random matrix model for elliptic curve $L$ functions of finite conductor. J. Phys. A 45 (2012).
\bibitem{diaconisevans2001} P. Diaconis and S. Evans. Linear functionals of eigenvalues of random matrices. Trans. Amer. Math. Soc. 353 (2001).
\bibitem{diaconiskemperman1996} P. Diaconis and J. Kemperman. Some new tools for Dirichlet priors. In Bayesian Statistics 5, Oxford University Press (1996).
\bibitem{diaconismostelleronishi1977} P. Diaconis, F. Mosteller and H. Onishi. Second order terms for the variances and covariances of the number of prime factors, including the square free case. J. Number Theory 9 (1977).
\bibitem{diaconisshahshahani1994} P. Diaconis and M. Shahshahani. On the eigenvalues of random matrices. J. Appl. Probab. 31A (1994).
\bibitem{diamantiskleban2009} N. Diamantis and P. Kleban. New percolation crossing formulas and second order modular forms. Commun. Number Theory Phys. 3 (2009).
\bibitem{diamondshurman2005} F. Diamond and J. Shurman. A First Course in Modular Forms. Graduate Texts in Mathematics 228, Springer (2005).
\bibitem{dickman1930} K. Dickman. On the frequency of numbers containing prime factors of a certain relative magnitude. Ark. Mat. Astr. Fys. 22A (1930).
\bibitem{dkm1985} D. Dummit, H. Kisilevsky and J. McKay. Multiplicative products of $\eta$ functions. Contemp. Math. 45 (1985).
\bibitem{dobner2021} A. Dobner. A proof of Newman's conjecture for the extended Selberg class. Acta Arith. 201 (2021).
\bibitem{doetsch1974} G. Doetsch. Introduction to the Theory and Application of the Laplace Transformation, translated by W. Nader. Springer (1974).
\bibitem{dokchitser2010} T. Dokchitser and V. Dokchitser. On the Birch Swinnerton Dyer quotients modulo squares. Ann. of Math. 172 (2010).
\bibitem{dokchitser2013} T. Dokchitser and V. Dokchitser. Identifying Frobenius elements in Galois groups. Algebra Number Theory 7 (2013).
\bibitem{donnellygrimmett1993} P. Donnelly and G. Grimmett. On the asymptotic distribution of large prime factors. J. London Math. Soc. 47 (1993).
\bibitem{doob1949} J. L. Doob. Application of the theory of martingales. In Le Calcul des Probabilit\'es et ses Applications, Colloq. Internat. CNRS 13 (1949).
\bibitem{doylesnell1984} P. G. Doyle and J. L. Snell. Random Walks and Electric Networks. Mathematical Association of America (1984).
\bibitem{dsrv2014} B. Duplantier, R. Rhodes, S. Sheffield and V. Vargas. Critical Gaussian multiplicative chaos: convergence of the derivative martingale. Ann. Probab. 42 (2014).
\bibitem{dsrv2014b} B. Duplantier, R. Rhodes, S. Sheffield and V. Vargas. Renormalization of critical Gaussian multiplicative chaos and KPZ relation. Comm. Math. Phys. 330 (2014).
\bibitem{dubedat2004} J. Dub\'edat. Critical percolation in annuli and SLE$_6$. Comm. Math. Phys. 245 (2004).
\bibitem{duke1988} W. Duke. Hyperbolic distribution problems and half integral weight Maass forms. Invent. Math. 92 (1988).
\bibitem{dymmckean1972} H. Dym and H. P. McKean. Fourier Series and Integrals. Academic Press (1972).
\bibitem{eaton1989} M. L. Eaton. Group Invariance Applications in Statistics. NSF CBMS Regional Conference Series in Probability and Statistics 1, IMS and ASA (1989).
\bibitem{eatongeorge2021} M. L. Eaton and E. I. George. Charles Stein and invariance: beginning with the Hunt Stein theorem. Ann. Statist. 49 (2021).
\bibitem{eichler1954} M. Eichler. Quatern\"are quadratische Formen und die Riemannsche Vermutung f\"ur die Kongruenzzetafunktion. Arch. Math. 5 (1954).
\bibitem{eichler1957} M. Eichler. Eine Verallgemeinerung der Abelschen Integrale. Math. Z. 67 (1957).
\bibitem{eichler1973} M. Eichler. The basis problem for modular forms and the traces of the Hecke operators. In Modular Functions of One Variable I, Lecture Notes in Math. 320, Springer (1973).
\bibitem{elboimgorodetsky2024} D. Elboim and O. Gorodetsky. Multiplicative arithmetic functions and the generalized Ewens measure. Israel J. Math. 262 (2024).
\bibitem{elkies1987} N. D. Elkies. The existence of infinitely many supersingular primes for every elliptic curve over $\mathbb{Q}$. Invent. Math. 89 (1987).
\bibitem{elkies1991} N. D. Elkies. Distribution of supersingular primes. Ast\'erisque 198, 199, 200 (1991).
\bibitem{erdoskac1940} P. Erd\H{o}s and M. Kac. The Gaussian law of errors in the theory of additive number theoretic functions. Amer. J. Math. 62 (1940).
\bibitem{erdosturan1948} P. Erd\H{o}s and P. Tur\'an. On a problem in the theory of uniform distribution I. Indag. Math. 10 (1948).
\bibitem{evelynlinfoot1931} C. J. A. Evelyn and E. H. Linfoot. On a problem in the additive theory of numbers IV. Ann. of Math. 32 (1931).
\bibitem{evw2016} J. S. Ellenberg, A. Venkatesh and C. Westerland. Homological stability for Hurwitz spaces and the Cohen Lenstra conjecture over function fields. Ann. of Math. 183 (2016).
\bibitem{ewens1972} W. J. Ewens. The sampling theory of selectively neutral alleles. Theoret. Population Biology 3 (1972).
\bibitem{faltings1983} G. Faltings. Endlichkeitss\"atze f\"ur abelsche Variet\"aten \"uber Zahlk\"orpern. Invent. Math. 73 (1983).
\bibitem{farguesscholze2026} L. Fargues and P. Scholze. Geometrization of the local Langlands correspondence. Ast\'erisque 466 (2026).
\bibitem{feller1968} W. Feller. An Introduction to Probability Theory and Its Applications, Volume I. Third edition, Wiley (1968).
\bibitem{feller1971} W. Feller. An Introduction to Probability Theory and Its Applications, Volume II, second edition. Wiley (1971).
\bibitem{ferguson1973} T. S. Ferguson. A Bayesian analysis of some nonparametric problems. Ann. Statist. 1 (1973).
\bibitem{fhk2012} Y. V. Fyodorov, G. A. Hiary and J. P. Keating. Freezing transition, characteristic polynomials of random matrices, and the Riemann zeta function. Phys. Rev. Lett. 108 (2012).
\bibitem{fill2009} J. A. Fill. The passage time distribution for a birth and death chain: strong stationary duality gives a first stochastic proof. J. Theoret. Probab. 22 (2009).
\bibitem{fisher1934} R. A. Fisher. Two new properties of mathematical likelihood. Proc. Roy. Soc. London A 144 (1934).
\bibitem{fisher1935} R. A. Fisher. The fiducial argument in statistical inference. Ann. Eugenics 6 (1935).
\bibitem{fishertippett1928} R. A. Fisher and L. H. C. Tippett. Limiting forms of the frequency distribution of the largest or smallest member of a sample. Proc. Cambridge Philos. Soc. 24 (1928).
\bibitem{fkrs2012} F. Fit\'e, K. Kedlaya, V. Rotger and A. Sutherland. Sato Tate distributions and Galois endomorphism modules in genus 2. Compos. Math. 148 (2012).
\bibitem{fks2021} F. Fit\'e, K. Kedlaya and A. Sutherland. Sato Tate groups of abelian threefolds: a preview of the classification. Contemp. Math. 770 (2021).
\bibitem{folland1995} G. B. Folland. A Course in Abstract Harmonic Analysis. CRC Press (1995).
\bibitem{fontainemazur1995} J. M. Fontaine and B. Mazur. Geometric Galois representations. In Elliptic Curves, Modular Forms and Fermat's Last Theorem (Hong Kong 1993), International Press (1995).
\bibitem{ford2022} K. Ford. Joint Poisson distribution of prime factors in sets. Math. Proc. Cambridge Philos. Soc. 173 (2022).
\bibitem{fordlucamoree2014} K. Ford, F. Luca and P. Moree. Values of the Euler phi function not divisible by a given odd prime, and the distribution of Euler Kronecker constants for cyclotomic fields. Math. Comp. 83 (2014).
\bibitem{fouvrykluners2007} \'E. Fouvry and J. Kl\"uners. On the 4 rank of class groups of quadratic number fields. Invent. Math. 167 (2007).
\bibitem{fouvrymichel2007} E. Fouvry and P. Michel. Sur le changement de signe des sommes de Kloosterman. Ann. of Math. 165 (2007).
\bibitem{fouvrymurty1996} E. Fouvry and M. R. Murty. On the distribution of supersingular primes. Canad. J. Math. 48 (1996).
\bibitem{fraser1968} D. A. S. Fraser. The Structure of Inference. Wiley (1968).
\bibitem{frenkel2007} E. Frenkel. Lectures on the Langlands program and conformal field theory. In Frontiers in Number Theory, Physics, and Geometry II, Springer (2007).
\bibitem{frey1986} G. Frey. Links between stable elliptic curves and certain Diophantine equations. Ann. Univ. Sarav. Ser. Math. 1 (1986).
\bibitem{friedman2008} J. Friedman. A proof of Alon's second eigenvalue conjecture and related problems. Mem. Amer. Math. Soc. 195 (2008).
\bibitem{friedmanwashington1989} E. Friedman and L. C. Washington. On the distribution of divisor class groups of curves over a finite field. In Th\'eorie des Nombres (Quebec 1987), de Gruyter (1989).
\bibitem{fyodorovkeating2014} Y. V. Fyodorov and J. P. Keating. Freezing transitions and extreme values: random matrix theory, $\zeta(1/2+it)$ and disordered landscapes. Philos. Trans. R. Soc. A 372 (2014).
\bibitem{gaitsgory2024} D. Gaitsgory, S. Raskin and collaborators. Proof of the geometric Langlands conjecture. Series of preprints (2024).
\bibitem{garunkstis2000} R. Garunk\v{s}tis. A note on the Riemann $\xi$ function (in Lithuanian). Liet. Mat. Rink. 40 (2000).
\bibitem{garunkstis2002} R. Garunk\v{s}tis. On a positivity property of the Riemann $\xi$ function. Lith. Math. J. 42 (2002).
\bibitem{gaudin1961} M. Gaudin. Sur la loi limite de l'espacement des valeurs propres d'une matrice al\'eatoire. Nuclear Phys. 25 (1961).
\bibitem{gauss1811} C. F. Gauss. Summatio quarumdam serierum singularium. Comment. Soc. Reg. Sci. Gotting. Recent. 1 (1811).
\bibitem{gelbart1984} S. Gelbart. An elementary introduction to the Langlands program. Bull. Amer. Math. Soc. 10 (1984).
\bibitem{gelbartjacquet1978} S. Gelbart and H. Jacquet. A relation between automorphic representations of GL(2) and GL(3). Ann. Sci. \'Ecole Norm. Sup. 11 (1978).
\bibitem{gelbartshahidi2001} S. Gelbart and F. Shahidi. Boundedness of automorphic $L$ functions in vertical strips. J. Amer. Math. Soc. 14 (2001).
\bibitem{ghk2007} S. M. Gonek, C. P. Hughes and J. P. Keating. A hybrid Euler Hadamard product for the Riemann zeta function. Duke Math. J. 136 (2007).
\bibitem{gmrr2021} O. Gorodetsky, K. Matom\"aki, M. Radziwi\l\l{} and B. Rodgers. On the variance of squarefree integers in short intervals and arithmetic progressions. Geom. Funct. Anal. 31 (2021).
\bibitem{gnedenko1943} B. Gnedenko. Sur la distribution limite du terme maximum d'une s\'erie al\'eatoire. Ann. of Math. 44 (1943).
\bibitem{godementjacquet1972} R. Godement and H. Jacquet. Zeta Functions of Simple Algebras. Lecture Notes in Math. 260, Springer (1972).
\bibitem{goldfeld1979} D. Goldfeld. Conjectures on elliptic curves over quadratic fields. In Number Theory, Carbondale 1979, Lecture Notes in Math. 751, Springer (1979).
\bibitem{goldfeld1982} D. Goldfeld. Sur les produits partiels eul\'eriens attach\'es aux courbes elliptiques. C. R. Acad. Sci. Paris 294 (1982).
\bibitem{good1958} I. J. Good. The interaction algorithm and practical Fourier analysis. J. Roy. Statist. Soc. Ser. B 20 (1958).
\bibitem{good1962} I. J. Good. Analogues of Poisson's summation formula. Amer. Math. Monthly 69 (1962).
\bibitem{goodgaskins1971} I. J. Good and R. A. Gaskins. Nonparametric roughness penalties for probability densities. Biometrika 58 (1971).
\bibitem{goodgaskins1980} I. J. Good and R. A. Gaskins. Density estimation and bump hunting by the penalized likelihood method exemplified by scattering and meteorite data. J. Amer. Statist. Assoc. 75 (1980).
\bibitem{gorodetskywong2025} O. Gorodetsky and M. D. Wong. On the limiting distribution of sums of random multiplicative functions. Preprint, arXiv 2508.12956 (2025).
\bibitem{grenandersz} U. Grenander and G. Szeg\H{o}. Toeplitz Forms and Their Applications. University of California Press (1958).
\bibitem{grommer1914} J. Grommer. Ganze transzendente Funktionen mit lauter reellen Nullstellen. J. Reine Angew. Math. 144 (1914).
\bibitem{gross1987} B. H. Gross. Heights and the special values of $L$ series. In Number Theory, CMS Conf. Proc. 7, Amer. Math. Soc. (1987).
\bibitem{grosszagier1986} B. Gross and D. Zagier. Heegner points and derivatives of L series. Invent. Math. 84 (1986).
\bibitem{gumbel1958} E. J. Gumbel. Statistics of Extremes. Columbia University Press (1958).
\bibitem{gutmanmohar1996} I. Gutman and B. Mohar. The quasi Wiener and the Kirchhoff indices coincide. J. Chem. Inf. Comput. Sci. 36 (1996).
\bibitem{haar1933} A. Haar. Der Massbegriff in der Theorie der kontinuierlichen Gruppen. Ann. of Math. 34 (1933).
\bibitem{hadamard1896} J. Hadamard. Sur la distribution des z\'eros de la fonction $\zeta(s)$ et ses cons\'equences arithm\'etiques. Bull. Soc. Math. France 24 (1896).
\bibitem{hafner1983} J. L. Hafner. Zeros on the critical line for Dirichlet series attached to certain cusp forms. Math. Ann. 264 (1983).
\bibitem{hall1938} P. Hall. A partition formula connected with Abelian groups. Comment. Math. Helv. 11 (1938).
\bibitem{hardy1914} G. H. Hardy. Sur les z\'eros de la fonction $\zeta(s)$ de Riemann. C. R. Acad. Sci. Paris 158 (1914).
\bibitem{harper2019} A. J. Harper. On the partition function of the Riemann zeta function, and the Fyodorov Hiary Keating conjecture. Preprint, arXiv 1906.05783 (2019).
\bibitem{harper2019b} A. J. Harper. Moments of random multiplicative functions II: high moments. Algebra Number Theory 13 (2019).
\bibitem{harper2020} A. J. Harper. Moments of random multiplicative functions I: low moments, better than squareroot cancellation, and critical multiplicative chaos. Forum Math. Pi 8 (2020).
\bibitem{harper2023} A. J. Harper. The typical size of character and zeta sums is $o(\sqrt x)$. Preprint, arXiv 2301.04390 (2023).
\bibitem{harris2009} M. Harris. Potential automorphy of odd dimensional symmetric powers of elliptic curves and applications. In Algebra, Arithmetic, and Geometry, Progress in Mathematics 270, Birkh\"auser (2009).
\bibitem{harristaylor2001} M. Harris and R. Taylor. The Geometry and Cohomology of Some Simple Shimura Varieties. Annals of Mathematics Studies 151, Princeton University Press (2001).
\bibitem{hartman1976} P. Hartman. Completely monotone families of solutions of $n$th order linear differential equations and infinitely divisible distributions. Ann. Scuola Norm. Sup. Pisa Cl. Sci. 3 (1976).
\bibitem{harvey2014} D. Harvey. Counting points on hyperelliptic curves in average polynomial time. Ann. of Math. 179 (2014).
\bibitem{harveysutherland2014} D. Harvey and A. V. Sutherland. Computing Hasse Witt matrices of hyperelliptic curves in average polynomial time. LMS J. Comput. Math. 17 (2014).
\bibitem{hasebesimonwang2018} T. Hasebe, T. Simon and M. Wang. Some properties of the free stable distributions. Ann. Inst. H. Poincar\'e Probab. Statist. 56 (2020).
\bibitem{hashimoto1989} K. Hashimoto. Zeta functions of finite graphs and representations of $p$ adic groups. Adv. Stud. Pure Math. 15 (1989).
\bibitem{heathbrown2004} D. R. Heath Brown. The average analytic rank of elliptic curves. Duke Math. J. 122 (2004).
\bibitem{hecke1918} E. Hecke. Eine neue Art von Zetafunktionen und ihre Beziehungen zur Verteilung der Primzahlen I, II. Math. Z. 1 (1918), 6 (1920).
\bibitem{hecke1920} E. Hecke. Eine neue Art von Zetafunktionen und ihre Beziehungen zur Verteilung der Primzahlen, II. Math. Z. 6 (1920).
\bibitem{hecke1936} E. Hecke. \"Uber die Bestimmung Dirichletscher Reihen durch ihre Funktionalgleichung. Math. Ann. 112 (1936).
\bibitem{hecke1937} E. Hecke. \"Uber Modulfunktionen und die Dirichletschen Reihen mit Eulerscher Produktentwicklung I, II. Math. Ann. 114 (1937).
\bibitem{hejhal1976} D. A. Hejhal. The Selberg Trace Formula for PSL(2,R), Volume 1. Lecture Notes in Math. 548, Springer (1976).
\bibitem{hejhal1983} D. A. Hejhal. The Selberg Trace Formula for PSL(2,R), Volume 2. Lecture Notes in Math. 1001, Springer (1983).
\bibitem{hejhal1992} D. A. Hejhal. Eigenvalues of the Laplacian for Hecke triangle groups. Mem. Amer. Math. Soc. 469 (1992).
\bibitem{hellegouarch1975} Y. Hellegouarch. Points d'ordre $2p^h$ sur les courbes elliptiques. Acta Arith. 26 (1975).
\bibitem{helson2010} H. Helson. Hankel forms. Studia Math. 198 (2010).
\bibitem{henniart2000} G. Henniart. Une preuve simple des conjectures de Langlands pour $\mathrm{GL}(n)$ sur un corps $p$ adique. Invent. Math. 139 (2000).
\bibitem{hewittsavage1955} E. Hewitt and L. J. Savage. Symmetric measures on Cartesian products. Trans. Amer. Math. Soc. 80 (1955).
\bibitem{hik2027} G. Hiary, S. Ireland and M. Kyi. A method for verifying the generalized Riemann hypothesis. Math. Comp. 96 (2027).
\bibitem{hinds1974} W. E. Hinds. Moments of complex random variables related to a certain class of characteristic functions. Sankhy\=a Ser. A 36 (1974).
\bibitem{hinkkanen1997} A. Hinkkanen. On functions of bounded type. Complex Variables Theory Appl. 34 (1997).
\bibitem{hko2001} C. P. Hughes, J. P. Keating and N. O'Connell. On the characteristic polynomial of a random unitary matrix. Comm. Math. Phys. 220 (2001).
\bibitem{hmy2024} J. Huang, T. McKenzie and H. T. Yau. Ramanujan property and edge universality of random regular graphs. Preprint, arXiv 2412.20263 (2024).
\bibitem{hoppe1984} F. M. Hoppe. P\'olya like urns and the Ewens' sampling formula. J. Math. Biol. 20 (1984).
\bibitem{horabuehler1966} R. B. Hora and R. J. Buehler. Fiducial theory and invariant estimation. Ann. Math. Statist. 37 (1966).
\bibitem{hsbt2010} M. Harris, N. Shepherd Barron and R. Taylor. A family of Calabi Yau varieties and potential automorphy. Ann. of Math. 171 (2010).
\bibitem{huber1974} P. J. Huber. Fisher information and spline interpolation. Ann. Statist. 2 (1974).
\bibitem{ichinoikeda2010} A. Ichino and T. Ikeda. On the periods of automorphic forms on special orthogonal groups and the Gross Prasad conjecture. Geom. Funct. Anal. 19 (2010).
\bibitem{ihara1966} Y. Ihara. On discrete subgroups of the two by two projective linear group over $p$ adic fields. J. Math. Soc. Japan 18 (1966).
\bibitem{ihara2006} Y. Ihara. On the Euler Kronecker constants of global fields and primes with small norms. In Algebraic Geometry and Number Theory, Progr. Math. 253, Birkh\"auser (2006).
\bibitem{ils2000} H. Iwaniec, W. Luo and P. Sarnak. Low lying zeros of families of $L$ functions. Publ. Math. IHES 91 (2000).
\bibitem{ingham1936} A. E. Ingham. Some trigonometrical inequalities with applications to the theory of series. Math. Z. 41 (1936).
\bibitem{irelandrosen1990} K. Ireland and M. Rosen. A Classical Introduction to Modern Number Theory, second edition. Graduate Texts in Mathematics 84, Springer (1990).
\bibitem{iwaniec1987} H. Iwaniec. Fourier coefficients of modular forms of half integral weight. Invent. Math. 87 (1987).
\bibitem{iwaniec2002} H. Iwaniec. Spectral Methods of Automorphic Forms. Second edition, Graduate Studies in Mathematics 53, American Mathematical Society (2002).
\bibitem{iwanieckowalski2004} H. Iwaniec and E. Kowalski. Analytic Number Theory. AMS Colloquium Publications 53 (2004).
\bibitem{jacquetlanglands1970} H. Jacquet and R. P. Langlands. Automorphic Forms on GL(2). Lecture Notes in Math. 114 (1970).
\bibitem{jacquetshalika1976} H. Jacquet and J. A. Shalika. A non vanishing theorem for zeta functions of $\mathrm{GL}_n$. Invent. Math. 38 (1976).
\bibitem{jacquetshalika1981} H. Jacquet and J. A. Shalika. On Euler products and the classification of automorphic representations I. Amer. J. Math. 103 (1981).
\bibitem{jamesroynetteyor2008} L. F. James, B. Roynette and M. Yor. Generalized gamma convolutions, Dirichlet means, Thorin measures, with explicit examples. Probab. Surveys 5 (2008).
\bibitem{jaodefeo2011} D. Jao and L. De Feo. Towards quantum resistant cryptosystems from supersingular elliptic curve isogenies. In Post Quantum Cryptography, PQCrypto 2011, Lecture Notes in Comput. Sci. 7071 (2011).
\bibitem{jeffreys1961} H. Jeffreys. Theory of Probability. Third edition, Oxford University Press (1961).
\bibitem{jessenwintner1935} B. Jessen and A. Wintner. Distribution functions and the Riemann zeta function. Trans. Amer. Math. Soc. 38 (1935).
\bibitem{johansson2017} C. Johansson. On the Sato Tate conjecture for non generic abelian surfaces. Trans. Amer. Math. Soc. 369 (2017).
\bibitem{jpss1981} H. Jacquet, I. I. Piatetski Shapiro and J. A. Shalika. Conducteur des repr\'esentations du groupe lin\'eaire. Math. Ann. 256 (1981).
\bibitem{jpss1983} H. Jacquet, I. I. Piatetski Shapiro and J. A. Shalika. Rankin Selberg convolutions. Amer. J. Math. 105 (1983).
\bibitem{jurek2010} Z. J. Jurek. Generalized L\'evy stochastic areas and selfdecomposability. Statist. Probab. Lett. 64 (2003).
\bibitem{jutila1981} M. Jutila. On the mean value of $L(1/2,\chi)$ for real characters. Analysis 1 (1981).
\bibitem{kahane1985} J. P. Kahane. Sur le chaos multiplicatif. Ann. Sci. Math. Qu\'ebec 9 (1985).
\bibitem{kamienny1992} S. Kamienny. Torsion points on elliptic curves and $q$ coefficients of modular forms. Invent. Math. 109 (1992).
\bibitem{kamiennynajman2012} S. Kamienny and F. Najman. Torsion groups of elliptic curves over quadratic fields. Acta Arith. 152 (2012).
\bibitem{karamata1931} J. Karamata. Neuer Beweis und Verallgemeinerung der Tauberschen S\"atze, welche die Laplacesche und Stieltjessche Transformation betreffen. J. Reine Angew. Math. 164 (1931).
\bibitem{karlin1968} S. Karlin. Total Positivity, Volume I. Stanford University Press (1968).
\bibitem{kassraftery1995} R. E. Kass and A. E. Raftery. Bayes factors. J. Amer. Statist. Assoc. 90 (1995).
\bibitem{katz1981} N. M. Katz. Galois properties of torsion points on abelian varieties. Invent. Math. 62 (1981).
\bibitem{katz1988} N. M. Katz. Gauss Sums, Kloosterman Sums, and Monodromy Groups. Annals of Mathematics Studies 116, Princeton University Press (1988).
\bibitem{katz1990} N. M. Katz. Exponential Sums and Differential Equations. Annals of Mathematics Studies 124, Princeton University Press (1990).
\bibitem{katzsarnak1999} N. M. Katz and P. Sarnak. Random Matrices, Frobenius Eigenvalues, and Monodromy. AMS Colloquium Publications 45 (1999).
\bibitem{katzsarnak1999b} N. M. Katz and P. Sarnak. Zeroes of zeta functions and symmetry. Bull. Amer. Math. Soc. 36 (1999).
\bibitem{keatingsnaith2000} J. Keating and N. Snaith. Random matrix theory and $\zeta(1/2+it)$. Comm. Math. Phys. 214 (2000).
\bibitem{keatingsnaith2000b} J. Keating and N. Snaith. Random matrix theory and L functions at $s=1/2$. Comm. Math. Phys. 214 (2000).
\bibitem{kedlayasutherland2009} K. Kedlaya and A. Sutherland. Hyperelliptic curves, L polynomials, and random matrices. Contemp. Math. 487 (2009).
\bibitem{keilson1971} J. Keilson. Log concavity and log convexity in passage time densities of diffusion and birth death processes. J. Appl. Probab. 8 (1971).
\bibitem{kemenysnell1960} J. G. Kemeny and J. L. Snell. Finite Markov Chains. Van Nostrand (1960).
\bibitem{kendall1957} D. G. Kendall. Some problems in the theory of dams. J. Roy. Statist. Soc. Ser. B 19 (1957).
\bibitem{kenkumomose1988} M. A. Kenku and F. Momose. Torsion points on elliptic curves defined over quadratic fields. Nagoya Math. J. 109 (1988).
\bibitem{kent1978} J. Kent. Some probabilistic properties of Bessel functions. Ann. Probab. 6 (1978).
\bibitem{kent1980} J. T. Kent. Eigenvalue expansions for diffusion hitting times. Z. Wahrsch. Verw. Gebiete 52 (1980).
\bibitem{kesten1959} H. Kesten. Symmetric random walks on groups. Trans. Amer. Math. Soc. 92 (1959).
\bibitem{kesten1959b} H. Kesten. Full Banach mean values on countable groups. Math. Scand. 7 (1959).
\bibitem{kharewintenberger2009} C. Khare and J. P. Wintenberger. Serre's modularity conjecture I. Invent. Math. 178 (2009).
\bibitem{kharewintenberger2009b} C. Khare and J. P. Wintenberger. Serre's modularity conjecture II. Invent. Math. 178 (2009).
\bibitem{khinchine1938} A. Ya. Khinchine. Limit Laws for Sums of Independent Random Variables (in Russian). ONTI, Moscow and Leningrad (1938).
\bibitem{kimmurty2023} S. Kim and M. R. Murty. From the Birch and Swinnerton Dyer conjecture to Nagao's conjecture. Math. Comp. 92 (2023).
\bibitem{kimsarnak2003} H. H. Kim. Functoriality for the exterior square of $\mathrm{GL}_4$ and the symmetric fourth of $\mathrm{GL}_2$, with appendix 1 by D. Ramakrishnan and appendix 2 by H. H. Kim and P. Sarnak. J. Amer. Math. Soc. 16 (2003).
\bibitem{kimshahidi2002} H. Kim and F. Shahidi. Functorial products for GL(2) $\times$ GL(3) and the symmetric cube for GL(2). Ann. of Math. 155 (2002).
\bibitem{kingman1975} J. F. C. Kingman. Random discrete distributions. J. Roy. Statist. Soc. Ser. B 37 (1975).
\bibitem{kingman1993} J. F. C. Kingman. Poisson Processes. Oxford University Press (1993).
\bibitem{kirchhoff1847} G. Kirchhoff. \"Uber die Aufl\"osung der Gleichungen, auf welche man bei der Untersuchung der linearen Vertheilung galvanischer Str\"ome gef\"uhrt wird. Ann. Phys. Chem. 72 (1847).
\bibitem{kisin2009} M. Kisin. Modularity of 2 adic Barsotti Tate representations. Invent. Math. 178 (2009).
\bibitem{kistler2015} N. Kistler. Derrida's random energy models: from spin glasses to the extremes of correlated random fields. In Correlated Random Systems: Five Different Methods, Lecture Notes in Math. 2143, Springer (2015).
\bibitem{klebanzagier2003} P. Kleban and D. Zagier. Crossing probabilities and modular forms. J. Stat. Phys. 113 (2003).
\bibitem{kleijnvandervaart2006} B. J. K. Kleijn and A. W. van der Vaart. Misspecification in infinite dimensional Bayesian statistics. Ann. Statist. 34 (2006).
\bibitem{kloosterman1926} H. D. Kloosterman. On the representation of numbers in the form $ax^2+by^2+cz^2+dt^2$. Acta Math. 49 (1926).
\bibitem{knapp1994} A. W. Knapp. Local Langlands correspondence: the Archimedean case. In Motives, Proc. Sympos. Pure Math. 55, Part 2, Amer. Math. Soc. (1994).
\bibitem{knuthtrabbpardo1976} D. E. Knuth and L. Trabb Pardo. Analysis of a simple factorization algorithm. Theoret. Comput. Sci. 3 (1976).
\bibitem{koblitz1988} N. Koblitz. Primality of the number of points on an elliptic curve over a finite field. Pacific J. Math. 131 (1988).
\bibitem{koehler2011} G. K\"ohler. Eta Products and Theta Series Identities. Springer Monographs in Mathematics, Springer (2011).
\bibitem{kolmogorov1933} A. N. Kolmogorov. Sulla determinazione empirica di una legge di distribuzione. Giorn. Ist. Ital. Attuari 4 (1933).
\bibitem{kolyvagin1988} V. A. Kolyvagin. Finiteness of $E(\mathbb{Q})$ and of the Shafarevich Tate group for a subclass of Weil curves. Izv. Akad. Nauk SSSR Ser. Mat. 52 (1988).
\bibitem{konstantopoulospatiesarkar2024} T. Konstantopoulos, P. Patie and R. Sarkar. A new class of solutions to the van Dantzig problem, the Lee Yang property, and the Riemann hypothesis. Ann. Inst. Fourier 74 (2024).
\bibitem{kowalskisawin2016} E. Kowalski and W. F. Sawin. Kloosterman paths and the shape of exponential sums. Compos. Math. 152 (2016).
\bibitem{kowalskisawin2017} E. Kowalski and W. F. Sawin. On the support of the Kloosterman paths. Preprint, arXiv 1709.05192 (2017).
\bibitem{kristiansen1994} G. K. Kristiansen. A proof of Steutel's conjecture. Ann. Probab. 22 (1994).
\bibitem{kubilius1955} J. Kubilius. On a problem in the $n$ dimensional analytic theory of numbers. Vilniaus Valst. Univ. Mokslo Darbai 4 (1955).
\bibitem{kullbackleibler1951} S. Kullback and R. A. Leibler. On information and sufficiency. Ann. Math. Statist. 22 (1951).
\bibitem{kuomurty2005} W. Kuo and M. R. Murty. On a conjecture of Birch and Swinnerton Dyer. Canad. J. Math. 57 (2005).
\bibitem{kuznetsov1980} N. V. Kuznetsov. Petersson's conjecture for cusp forms of weight zero and Linnik's conjecture. Sums of Kloosterman sums. Mat. Sb. 111 (1980); English translation Math. USSR Sb. 39 (1981).
\bibitem{kyprianou2014} A. E. Kyprianou. Fluctuations of L\'evy Processes with Applications. Second edition, Universitext, Springer (2014).
\bibitem{lagarias1999} J. C. Lagarias. On a positivity property of the Riemann $\xi$ function. Acta Arith. 89 (1999).
\bibitem{lagariasodlyzko1977} J. C. Lagarias and A. M. Odlyzko. Effective versions of the Chebotarev density theorem. In Algebraic Number Fields, Academic Press (1977).
\bibitem{lalleysellke1987} S. P. Lalley and T. Sellke. A conditional limit theorem for the frontier of a branching Brownian motion. Ann. Probab. 15 (1987).
\bibitem{lamzouri2011} Y. Lamzouri. On the distribution of extreme values of zeta and $L$ functions in the strip. Int. Math. Res. Not. IMRN (2011).
\bibitem{lamzouri2020} Y. Lamzouri. On the distribution of the maximum of cubic exponential sums. J. Inst. Math. Jussieu 19 (2020).
\bibitem{landau1905} E. Landau. \"Uber einen Satz von Tschebyschef. Math. Ann. 61 (1905).
\bibitem{landau1908} E. Landau. \"Uber die Einteilung der positiven ganzen Zahlen in vier Klassen nach der Mindestzahl der zu ihrer additiven Zusammensetzung erforderlichen Quadrate. Arch. Math. Phys. 13 (1908).
\bibitem{landsberg1893} G. Landsberg. Zur Theorie der Gaussschen Summen und der linearen Transformation der Thetafunctionen. J. Reine Angew. Math. 111 (1893).
\bibitem{lang1976} S. Lang. Introduction to Modular Forms, with appendices by D. Zagier and W. Feit. Grundlehren 222, Springer (1976).
\bibitem{lang1995} S. Lang. Some history of the Shimura Taniyama conjecture. Notices Amer. Math. Soc. 42 (1995).
\bibitem{langlands1967} R. P. Langlands. Letter to Andr\'e Weil, January 1967. The Work of Robert Langlands, Institute for Advanced Study (1967).
\bibitem{langlands1970} R. P. Langlands. Problems in the theory of automorphic forms. In Lectures in Modern Analysis and Applications III, Lecture Notes in Math. 170, Springer (1970).
\bibitem{langlands1980} R. P. Langlands. Base Change for GL(2). Annals of Mathematics Studies 96, Princeton University Press (1980).
\bibitem{langlands1989} R. P. Langlands. On the classification of irreducible representations of real algebraic groups. In Representation Theory and Harmonic Analysis on Semisimple Lie Groups, Math. Surveys Monogr. 31, Amer. Math. Soc. (1989).
\bibitem{langtrotter1976} S. Lang and H. Trotter. Frobenius Distributions in $\mathrm{GL}_2$ Extensions. Lecture Notes in Math. 504, Springer (1976).
\bibitem{languascozaccagnini2007} A. Languasco and A. Zaccagnini. A note on Mertens' formula for arithmetic progressions. J. Number Theory 127 (2007).
\bibitem{laurincikasmatsumoto2001} A. Laurin\v{c}ikas and K. Matsumoto. The universality of zeta functions attached to certain cusp forms. Acta Arith. 98 (2001).
\bibitem{lavine1992} M. Lavine. Some aspects of P\'olya tree distributions for statistical modelling. Ann. Statist. 20 (1992).
\bibitem{lazarevic1966} I. Lazarevi\'c. Neke nejednakosti sa hiperboli\v{c}kim funkcijama. Univ. Beograd. Publ. Elektrotehn. Fak. Ser. Mat. Fiz. 170 (1966).
\bibitem{leadbetter1983} M. R. Leadbetter, G. Lindgren and H. Rootz\'en. Extremes and Related Properties of Random Sequences and Processes. Springer (1983).
\bibitem{leeyang1952} T. D. Lee and C. N. Yang. Statistical theory of equations of state and phase transitions II. Lattice gas and Ising model. Phys. Rev. 87 (1952).
\bibitem{lehmer1947} D. H. Lehmer. The vanishing of Ramanujan's function $\tau(n)$. Duke Math. J. 14 (1947).
\bibitem{levin1964} B. Ja. Levin. Distribution of Zeros of Entire Functions. Translations of Mathematical Monographs 5, American Mathematical Society (1964).
\bibitem{levinsonmontgomery1974} N. Levinson and H. L. Montgomery. Zeros of the derivatives of the Riemann zeta function. Acta Math. 133 (1974).
\bibitem{li1997} X. J. Li. The positivity of a sequence of numbers and the Riemann hypothesis. J. Number Theory 65 (1997).
\bibitem{li2019} W. C. W. Li. Zeta and $L$ Functions in Number Theory and Combinatorics. CBMS Regional Conf. Ser. Math. 129, Amer. Math. Soc. (2019).
\bibitem{lindnerpansato2018} A. Lindner, L. Pan and K. Sato. On quasi infinitely divisible distributions. Trans. Amer. Math. Soc. 370 (2018).
\bibitem{linhu2001} G. D. Lin and C. Y. Hu. The Riemann zeta distribution. Bernoulli 7 (2001).
\bibitem{linnik1962} Yu. V. Linnik. Additive problems and eigenvalues of the modular operators. Proc. Internat. Congr. Mathematicians, Stockholm (1962).
\bibitem{linnik1968} Yu. V. Linnik. Ergodic Properties of Algebraic Fields. Springer (1968).
\bibitem{liu2016} H. Q. Liu. On the distribution of squarefree numbers. J. Number Theory 159 (2016).
\bibitem{lloyd1984} S. P. Lloyd. Ordered prime divisors of a random integer. Ann. Probab. 12 (1984).
\bibitem{lmfdb} The LMFDB Collaboration. The L functions and modular forms database. https://www.lmfdb.org.
\bibitem{louboutinmunsch2013} S. R. Louboutin and M. Munsch. On positive real zeros of theta and $L$ functions associated with real, even and primitive characters. Publ. Math. Debrecen 83 (2013).
\bibitem{lpps1992} R. P. Langlands, C. Pichet, P. Pouliot and Y. Saint Aubin. On the universality of crossing probabilities in two dimensional percolation. J. Stat. Phys. 67 (1992).
\bibitem{lps1988} A. Lubotzky, R. Phillips and P. Sarnak. Ramanujan graphs. Combinatorica 8 (1988).
\bibitem{lps1994} R. P. Langlands, P. Pouliot and Y. Saint Aubin. Conformal invariance in two dimensional percolation. Bull. Amer. Math. Soc. 30 (1994).
\bibitem{lrs1995} W. Luo, Z. Rudnick and P. Sarnak. On Selberg's eigenvalue conjecture. Geom. Funct. Anal. 5 (1995).
\bibitem{lukacs1968} E. Lukacs. Contributions to a problem of D. van Dantzig. Theory Probab. Appl. 13 (1968).
\bibitem{lukacs1970} E. Lukacs. Characteristic Functions. Second edition, Griffin (1970).
\bibitem{lyons2005} R. Lyons. Asymptotic enumeration of spanning trees. Combin. Probab. Comput. 14 (2005).
\bibitem{maass1949} H. Maass. \"Uber eine neue Art von nichtanalytischen automorphen Funktionen und die Bestimmung Dirichletscher Reihen durch Funktionalgleichungen. Math. Ann. 121 (1949).
\bibitem{macdonald1971} I. G. Macdonald. Spherical Functions on a Group of $p$ adic Type. Ramanujan Institute, Madras (1971).
\bibitem{malle2008} G. Malle. Cohen Lenstra heuristic and roots of unity. J. Number Theory 128 (2008).
\bibitem{margulis1988} G. A. Margulis. Explicit group theoretic constructions of combinatorial schemes and their applications in the construction of expanders and concentrators. Problems Inform. Transmission 24 (1988).
\bibitem{martin1996} Y. Martin. Multiplicative $\eta$ quotients. Trans. Amer. Math. Soc. 348 (1996).
\bibitem{mathstackexchange2016} Dirichlet series associated to squared M\"obius. Mathematics Stack Exchange, question 2063166 (2016).
\bibitem{mazur1977} B. Mazur. Modular curves and the Eisenstein ideal. Publ. Math. IHES 47 (1977).
\bibitem{mazur1978} B. Mazur. Rational isogenies of prime degree. Invent. Math. 44 (1978).
\bibitem{mazur2008} B. Mazur. Finding meaning in error terms. Bull. Amer. Math. Soc. 45 (2008).
\bibitem{mckay1981} B. D. McKay. The expected eigenvalue distribution of a large regular graph. Linear Algebra Appl. 40 (1981).
\bibitem{mckay1983} B. D. McKay. Spanning trees in regular graphs. European J. Combin. 4 (1983).
\bibitem{mckean1972} H. P. McKean. Selberg's trace formula as applied to a compact Riemann surface. Comm. Pure Appl. Math. 25 (1972).
\bibitem{meckes2019} E. S. Meckes. The Random Matrix Theory of the Classical Compact Groups. Cambridge Tracts in Mathematics 218, Cambridge University Press (2019).
\bibitem{mehta2004} M. L. Mehta. Random Matrices. Third edition, Elsevier (2004).
\bibitem{meng2017} X. Meng. The distribution of $k$ free numbers and the derivative of the Riemann zeta function. Math. Proc. Cambridge Philos. Soc. 162 (2017).
\bibitem{mertens1874} F. Mertens. Ein Beitrag zur analytischen Zahlentheorie. J. Reine Angew. Math. 78 (1874).
\bibitem{mestre1986} J. F. Mestre. La m\'ethode des graphes. Exemples et applications. In Proceedings of the International Conference on Class Numbers and Fundamental Units of Algebraic Number Fields, Katata (1986).
\bibitem{miller2004} S. J. Miller. One and two level densities for rational families of elliptic curves: evidence for the underlying group symmetries. Compos. Math. 140 (2004).
\bibitem{miller2006} S. J. Miller. Investigations of zeros near the central point of elliptic curve L functions. Experiment. Math. 15 (2006).
\bibitem{moeglinwaldspurger1989} C. Moeglin and J. L. Waldspurger. Le spectre r\'esiduel de $\mathrm{GL}(n)$. Ann. Sci. \'Ecole Norm. Sup. 22 (1989).
\bibitem{monsky1996} P. Monsky. Generalizing the Birch Stephens theorem I. Modular curves. Math. Z. 221 (1996).
\bibitem{montgomery1973} H. Montgomery. The pair correlation of zeros of the zeta function. Proc. Sympos. Pure Math. 24 (1973).
\bibitem{montgomeryvaughan1981} H. L. Montgomery and R. C. Vaughan. The distribution of squarefree numbers. In Recent Progress in Analytic Number Theory, Vol. 1, Academic Press (1981).
\bibitem{montgomeryvaughan2007} H. L. Montgomery and R. C. Vaughan. Multiplicative Number Theory I: Classical Theory. Cambridge University Press (2007).
\bibitem{mop2022} S. Mohanty, R. O'Donnell and P. Paredes. Explicit near Ramanujan graphs of every degree. SIAM J. Comput. 51 (2022).
\bibitem{mordell1917} L. J. Mordell. On Mr. Ramanujan's empirical expansions of modular functions. Proc. Cambridge Philos. Soc. 19 (1917).
\bibitem{morgenstern1994} M. Morgenstern. Existence and explicit constructions of $q+1$ regular Ramanujan graphs for every prime power $q$. J. Combin. Theory Ser. B 62 (1994).
\bibitem{mot2021} M. J. Mossinghoff, T. Oliveira e Silva and T. S. Trudgian. The distribution of $k$ free numbers. Math. Comp. 90 (2021).
\bibitem{mss2015} A. W. Marcus, D. A. Spielman and N. Srivastava. Interlacing families I: bipartite Ramanujan graphs of all degrees. Ann. of Math. 182 (2015).
\bibitem{mss2018} A. W. Marcus, D. A. Spielman and N. Srivastava. Interlacing families IV: bipartite Ramanujan graphs of all sizes. SIAM J. Comput. 47 (2018).
\bibitem{muniershnidman2023} N. Munier and A. Shnidman. Sandpile groups of supersingular isogeny graphs. J. Th\'eor. Nombres Bordeaux 35 (2023).
\bibitem{murty1985} V. K. Murty. Explicit formulae and the Lang Trotter conjecture. Rocky Mountain J. Math. 15 (1985).
\bibitem{murty1997} V. K. Murty. Modular forms and the Chebotarev density theorem II. In Analytic Number Theory (Kyoto, 1996), London Math. Soc. Lecture Note Ser. 247 (1997).
\bibitem{murtymurtysaradha1988} M. R. Murty, V. K. Murty and N. Saradha. Modular forms and the Chebotarev density theorem. Amer. J. Math. 110 (1988).
\bibitem{murtysinha2009} M. R. Murty and K. Sinha. Effective equidistribution of eigenvalues of Hecke operators. J. Number Theory 129 (2009).
\bibitem{nagoshi2006} H. Nagoshi. Distribution of Hecke eigenvalues. Proc. Amer. Math. Soc. 134 (2006).
\bibitem{najnudel2018} J. Najnudel. On the extreme values of the Riemann zeta function on random intervals of the critical line. Probab. Theory Related Fields 172 (2018).
\bibitem{neukirch1999} J. Neukirch. Algebraic Number Theory. Grundlehren 322, Springer (1999).
\bibitem{newman1974} C. M. Newman. Zeros of the partition function for generalized Ising systems. Comm. Pure Appl. Math. 27 (1974).
\bibitem{newman1976} C. M. Newman. Fourier transforms with only real zeros. Proc. Amer. Math. Soc. 61 (1976).
\bibitem{newmanwu2020} C. M. Newman and W. Wu. Constants of de Bruijn Newman type in analytic number theory and statistical physics. Bull. Amer. Math. Soc. 57 (2020).
\bibitem{newtonthorne2021} J. Newton and J. A. Thorne. Symmetric power functoriality for holomorphic modular forms. Publ. Math. IHES 134 (2021).
\bibitem{newtonthorne2021b} J. Newton and J. A. Thorne. Symmetric power functoriality for holomorphic modular forms II. Publ. Math. IHES 134 (2021).
\bibitem{ngo2010} B. C. Ng\^o. Le lemme fondamental pour les alg\`ebres de Lie. Publ. Math. IHES 111 (2010).
\bibitem{niederreiter1991} H. Niederreiter. The distribution of values of Kloosterman sums. Arch. Math. 56 (1991).
\bibitem{nilli1991} A. Nilli. On the second eigenvalue of a graph. Discrete Math. 91 (1991).
\bibitem{nsw2020} M. Nikula, E. Saksman and C. Webb. Multiplicative chaos and the characteristic polynomial of the CUE: the $L^1$ phase. Trans. Amer. Math. Soc. 373 (2020).
\bibitem{odlyzko1987} A. M. Odlyzko. On the distribution of spacings between zeros of the zeta function. Math. Comp. 48 (1987).
\bibitem{odlyzkoschonhage1988} A. M. Odlyzko and A. Sch\"onhage. Fast algorithms for multiple evaluations of the Riemann zeta function. Trans. Amer. Math. Soc. 309 (1988).
\bibitem{odoni1975} R. W. K. Odoni. On the norms of algebraic integers. Mathematika 22 (1975).
\bibitem{ogg1975} A. P. Ogg. Diophantine equations and modular forms. Bull. Amer. Math. Soc. 81 (1975).
\bibitem{ostrovskii1970} I. V. Ostrovskii. A certain class of characteristic functions. Trudy Mat. Inst. Steklov. 111 (1970).
\bibitem{paquettezeitouni2018} E. Paquette and O. Zeitouni. The maximum of the CUE field. Int. Math. Res. Not. IMRN (2018).
\bibitem{paquettezeitouni2022} E. Paquette and O. Zeitouni. The extremal landscape for the C$\beta$E ensemble. Forum Math. Sigma 13 (2025).
\bibitem{pari2023} The PARI Group. PARI/GP version 2.15. Universit\'e de Bordeaux (2023).
\bibitem{patiesavov2018} P. Patie and M. Savov. Bernstein gamma functions and exponential functionals of L\'evy processes. Electron. J. Probab. 23 (2018).
\bibitem{pawlewicz2011} J. Pawlewicz. Counting square free numbers. Preprint, arXiv 1107.4890 (2011).
\bibitem{perron1908} O. Perron. Zur Theorie der Dirichletschen Reihen. J. Reine Angew. Math. 134 (1908).
\bibitem{peterweyl1927} F. Peter and H. Weyl. Die Vollst\"andigkeit der primitiven Darstellungen einer geschlossenen kontinuierlichen Gruppe. Math. Ann. 97 (1927).
\bibitem{pitman2006} J. Pitman. Combinatorial Stochastic Processes. Lecture Notes in Mathematics 1875, Springer (2006).
\bibitem{pizer1990} A. K. Pizer. Ramanujan graphs and Hecke operators. Bull. Amer. Math. Soc. 23 (1990).
\bibitem{platt2016} D. J. Platt. Numerical computations concerning the GRH. Math. Comp. 85 (2016).
\bibitem{platttrudgian2021} D. Platt and T. Trudgian. The Riemann hypothesis is true up to $3\cdot10^{12}$. Bull. London Math. Soc. 53 (2021).
\bibitem{polson2017} N. G. Polson. On Hilbert's 8th problem. arXiv 1708.02653 (2017, version 26 2026).
\bibitem{polson2018vdp} N. G. Polson. Riemann, Thorin, van Dantzig pairs, Wald couples and Hadamard factorisation. arXiv 1804.10043 (2018, version 9 2026).
\bibitem{polson2026thorin} N. G. Polson. Thorin positivity, GGC factorization and the Riemann Xi function. Preprint, SSRN 6932859 (2026).
\bibitem{polya1926} G. P\'olya. Bemerkung \"uber die Integraldarstellung der Riemannschen $\xi$ Funktion. Acta Math. 48 (1926).
\bibitem{polya1927} G. P\'olya. \"Uber trigonometrische Integrale mit nur reellen Nullstellen. J. Reine Angew. Math. 158 (1927).
\bibitem{polymath2019} D. H. J. Polymath. Effective approximation of heat flow evolution of the Riemann $\xi$ function, and a new upper bound for the de Bruijn Newman constant. Res. Math. Sci. 6 (2019).
\bibitem{poonenrains2012} B. Poonen and E. Rains. Random maximal isotropic subspaces and Selmer groups. J. Amer. Math. Soc. 25 (2012).
\bibitem{ppvw2019} J. Park, B. Poonen, J. Voight and M. M. Wood. A heuristic for boundedness of ranks of elliptic curves. J. Eur. Math. Soc. 21 (2019).
\bibitem{rademachergrosswald1972} H. Rademacher and E. Grosswald. Dedekind Sums. Carus Mathematical Monographs 16, Mathematical Association of America (1972).
\bibitem{radziwillsound2017} M. Radziwi\l\l\ and K. Soundararajan. Selberg's central limit theorem for $\log|\zeta(1/2+it)|$. Enseign. Math. 63 (2017).
\bibitem{ramakrishna2002} R. Ramakrishna. Deforming Galois representations and the conjectures of Serre and Fontaine Mazur. Ann. of Math. 156 (2002).
\bibitem{ramakrishnan1994} D. Ramakrishnan. A refinement of the strong multiplicity one theorem for $\mathrm{GL}(2)$. Invent. Math. 116 (1994).
\bibitem{ramakrishnan2000} D. Ramakrishnan. Modularity of the Rankin Selberg L series, and multiplicity one for SL(2). Ann. of Math. 152 (2000).
\bibitem{ramanujan1916} S. Ramanujan. On certain arithmetical functions. Trans. Cambridge Philos. Soc. 22 (1916).
\bibitem{reid1988} N. Reid. Saddlepoint methods and statistical inference. Statist. Sci. 3 (1988).
\bibitem{renyi1953} A. R\'enyi. On the theory of order statistics. Acta Math. Acad. Sci. Hungar. 4 (1953).
\bibitem{resnick1987} S. I. Resnick. Extreme Values, Regular Variation, and Point Processes. Springer (1987).
\bibitem{rezvyakova2010} I. S. Rezvyakova. On the critical line zeros of L functions attached to automorphic cusp forms. Math. Notes 88 (2010).
\bibitem{rhodesvargas2014} R. Rhodes and V. Vargas. Gaussian multiplicative chaos and applications: a review. Probab. Surv. 11 (2014).
\bibitem{ribet1977} K. A. Ribet. Galois representations attached to eigenforms with Nebentypus. In Modular Functions of One Variable V, Lecture Notes in Math. 601, Springer (1977).
\bibitem{ribet1984} K. A. Ribet. Congruence relations between modular forms. In Proc. International Congress of Mathematicians, Warsaw 1983, Vol. 1 (1984).
\bibitem{ribet1990} K. A. Ribet. On modular representations of $\Gal(\overline\Q/\Q)$ arising from modular forms. Invent. Math. 100 (1990).
\bibitem{ribet1990b} K. A. Ribet. Raising the levels of modular representations. In S\'eminaire de Th\'eorie des Nombres, Paris 1987/88, Progr. Math. 81, Birkh\"auser (1990).
\bibitem{ricottaroyer2018} G. Ricotta and E. Royer. Kloosterman paths of prime powers moduli. Comment. Math. Helv. 93 (2018).
\bibitem{riemann1859} B. Riemann. \"Uber die Anzahl der Primzahlen unter einer gegebenen Gr\"osse. Monatsber. Berlin. Akad. (1859).
\bibitem{robbins1956} H. Robbins. An empirical Bayes approach to statistics. Proc. Third Berkeley Symp. Math. Statist. Probab. 1 (1956).
\bibitem{rodgerstao2020} B. Rodgers and T. Tao. The de Bruijn Newman constant is nonnegative. Forum Math. Pi 8 (2020).
\bibitem{rousethorner2017} J. Rouse and J. Thorner. The explicit Sato Tate conjecture and densities pertaining to Lehmer type questions. Trans. Amer. Math. Soc. 369 (2017).
\bibitem{roynetteyor2005} B. Roynette and M. Yor. Couples de Wald ind\'efiniment divisibles. Exemples li\'es \`a la fonction gamma d'Euler et \`a la fonction zeta de Riemann. Ann. Inst. Fourier 55 (2005).
\bibitem{rszb2022} J. Rouse, A. Sutherland and D. Zureick Brown. $\ell$ adic images of Galois for elliptic curves over $\mathbb{Q}$. Forum Math. Sigma 10 (2022).
\bibitem{rubinstein2005} M. Rubinstein. Computational methods and experiments in analytic number theory. In Recent Perspectives in Random Matrix Theory and Number Theory, London Math. Soc. Lecture Note Ser. 322, Cambridge University Press (2005).
\bibitem{rubinsteinsarnak1994} M. Rubinstein and P. Sarnak. Chebyshev's bias. Experiment. Math. 3 (1994).
\bibitem{rudnicksarnak1996} Z. Rudnick and P. Sarnak. Zeros of principal L functions and random matrix theory. Duke Math. J. 81 (1996).
\bibitem{rumely1993} R. Rumely. Numerical computations concerning the ERH. Math. Comp. 61 (1993).
\bibitem{saintaubin2021} Y. Saint Aubin. L'invariance conforme et l'universalit\'e au point critique des mod\`eles bidimensionnels. In The Genesis of the Langlands Program, Cambridge University Press (2021).
\bibitem{saksmanwebb2020} E. Saksman and C. Webb. The Riemann zeta function and Gaussian multiplicative chaos: statistics on the critical line. Ann. Probab. 48 (2020).
\bibitem{salminenvignat2024} P. Salminen and C. Vignat. Probabilistic aspects of Jacobi theta functions. Math. Scand. 130 (2024).
\bibitem{sarnak1982} P. Sarnak. Class numbers of indefinite binary quadratic forms. J. Number Theory 15 (1982).
\bibitem{sarnak1987} P. Sarnak. Statistical properties of eigenvalues of the Hecke operators. In Analytic Number Theory and Diophantine Problems, Progr. Math. 70, Birkh\"auser (1987).
\bibitem{sarnak1987b} P. Sarnak. Determinants of Laplacians. Comm. Math. Phys. 110 (1987).
\bibitem{sarnak2003} P. Sarnak. Spectra of hyperbolic surfaces. Bull. Amer. Math. Soc. 40 (2003).
\bibitem{sarnak2004} P. Sarnak. Nonvanishing of $L$ functions on $\Re(s)=1$. In Contributions to Automorphic Forms, Geometry, and Number Theory, Johns Hopkins University Press (2004).
\bibitem{sarnak2007} P. Sarnak. Letter to B. Mazur on Chebyshev's bias for $\tau(p)$ (2007).
\bibitem{sarnakreciprocal2007} P. Sarnak. Reciprocal geodesics. In Analytic Number Theory, Clay Math. Proc. 7, American Mathematical Society (2007).
\bibitem{sarnakxue1991} P. Sarnak and X. Xue. Bounds for multiplicities of automorphic representations. Duke Math. J. 64 (1991).
\bibitem{satake1963} I. Satake. Theory of spherical functions on reductive algebraic groups over $p$ adic fields. Publ. Math. IHES 18 (1963).
\bibitem{sathe1953} L. G. Sathe. On a problem of Hardy on the distribution of integers having a given number of prime factors. I, II. J. Indian Math. Soc. 17 (1953).
\bibitem{sato1999} K. Sato. L\'evy Processes and Infinitely Divisible Distributions. Cambridge University Press (1999).
\bibitem{schaar1850} M. Schaar. M\'emoire sur la th\'eorie des r\'esidus quadratiques. Mem. Acad. Roy. Sci. Lettres Beaux Arts Belgique 24 (1850).
\bibitem{schillingsongvondracek2012} R. L. Schilling, R. Song and Z. Vondra\v{c}ek. Bernstein Functions: Theory and Applications. Second edition, De Gruyter Studies in Mathematics 37, de Gruyter (2012).
\bibitem{schoenberg1951} I. J. Schoenberg. On P\'olya frequency functions I. The totally positive functions and their Laplace transforms. J. Analyse Math. 1 (1951).
\bibitem{scholze2013} P. Scholze. The local Langlands correspondence for GL(n) over $p$ adic fields. Invent. Math. 192 (2013).
\bibitem{scholze2015} P. Scholze. On torsion in the cohomology of locally symmetric varieties. Ann. of Math. 182 (2015).
\bibitem{schoof1985} R. Schoof. Elliptic curves over finite fields and the computation of square roots mod $p$. Math. Comp. 44 (1985).
\bibitem{schur1921} I. Schur. \"Uber die Gaussschen Summen. Nachr. Ges. Wiss. G\"ottingen (1921).
\bibitem{selberg1942} A. Selberg. On the zeros of Riemann's zeta function. Skr. Norske Vid. Akad. Oslo I (1942).
\bibitem{selberg1946} A. Selberg. Contributions to the theory of the Riemann zeta function. Arch. Math. Naturvid. 48 (1946).
\bibitem{selberg1954} A. Selberg. Note on a paper by L. G. Sathe. J. Indian Math. Soc. 18 (1954).
\bibitem{selberg1956} A. Selberg. Harmonic analysis and discontinuous groups in weakly symmetric Riemannian spaces with applications to Dirichlet series. J. Indian Math. Soc. 20 (1956).
\bibitem{selberg1965} A. Selberg. On the estimation of Fourier coefficients of modular forms. Proc. Sympos. Pure Math. 8 (1965).
\bibitem{selberg1992} A. Selberg. Old and new conjectures and results about a class of Dirichlet series. In Proceedings of the Amalfi Conference on Analytic Number Theory, Universit\`a di Salerno (1992).
\bibitem{serre1968} J. P. Serre. Abelian $\ell$ adic Representations and Elliptic Curves. Benjamin (1968).
\bibitem{serre1968b} J. P. Serre. Une interpr\'etation des congruences relatives \`a la fonction $\tau$ de Ramanujan. S\'eminaire Delange Pisot Poitou 9, expos\'e 14 (1968).
\bibitem{serre1970} J. P. Serre. Facteurs locaux des fonctions z\^eta des vari\'et\'es alg\'ebriques (d\'efinitions et conjectures). S\'eminaire Delange Pisot Poitou 11, expos\'e 19 (1970).
\bibitem{serre1972} J. P. Serre. Propri\'et\'es galoisiennes des points d'ordre fini des courbes elliptiques. Invent. Math. 15 (1972).
\bibitem{serre1973} J. P. Serre. A Course in Arithmetic. Graduate Texts in Mathematics 7, Springer (1973).
\bibitem{serre1976} J. P. Serre. Divisibilit\'e de certaines fonctions arithm\'etiques. Enseign. Math. 22 (1976).
\bibitem{serre1980} J. P. Serre. Trees. Springer (1980).
\bibitem{serre1981} J. P. Serre. Quelques applications du th\'eor\`eme de densit\'e de Chebotarev. Publ. Math. IHES 54 (1981).
\bibitem{serre1987} J. P. Serre. Sur les repr\'esentations modulaires de degr\'e 2 de $\Gal(\overline\Q/\Q)$. Duke Math. J. 54 (1987).
\bibitem{serre1997} J. P. Serre. R\'epartition asymptotique des valeurs propres de l'op\'erateur de Hecke $T_p$. J. Amer. Math. Soc. 10 (1997).
\bibitem{serre2012} J. P. Serre. Lectures on $N_X(p)$. CRC Press (2012).
\bibitem{shahidi1981} F. Shahidi. On certain $L$ functions. Amer. J. Math. 103 (1981).
\bibitem{shen2021} Q. Shen. The fourth moment of quadratic Dirichlet $L$ functions. Math. Z. 298 (2021).
\bibitem{shenstucky2024} Q. Shen and J. Stucky. The fourth moment of quadratic Dirichlet $L$ functions II. Preprint, arXiv 2402.01497 (2024).
\bibitem{shepplloyd1966} L. A. Shepp and S. P. Lloyd. Ordered cycle lengths in a random permutation. Trans. Amer. Math. Soc. 121 (1966).
\bibitem{shimura1971} G. Shimura. Introduction to the Arithmetic Theory of Automorphic Functions. Princeton University Press (1971).
\bibitem{shin2012} S. W. Shin. Automorphic Plancherel density theorem. Israel J. Math. 192 (2012).
\bibitem{shintemplier2016} S. W. Shin and N. Templier. Sato Tate theorem for families and low lying zeros of automorphic $L$ functions. Invent. Math. 203 (2016).
\bibitem{siegel1944} C. L. Siegel. The average measure of quadratic forms with given determinant and signature. Ann. of Math. 45 (1944).
\bibitem{siegel1945} C. L. Siegel. A mean value theorem in geometry of numbers. Ann. of Math. 46 (1945).
\bibitem{silverman1994} J. H. Silverman. Advanced Topics in the Arithmetic of Elliptic Curves. Graduate Texts in Mathematics 151, Springer (1994).
\bibitem{skinnerurban2014} C. Skinner and E. Urban. The Iwasawa main conjectures for $\mathrm{GL}_2$. Invent. Math. 195 (2014).
\bibitem{smirnov2001} S. Smirnov. Critical percolation in the plane: conformal invariance, Cardy's formula, scaling limits. C. R. Acad. Sci. Paris S\'er. I Math. 333 (2001).
\bibitem{smith2017} A. Smith. $2^\infty$ Selmer groups, $2^\infty$ class groups, and Goldfeld's conjecture. Preprint, arXiv 1702.02325 (2017).
\bibitem{soundararajan2000} K. Soundararajan. Nonvanishing of quadratic Dirichlet $L$ functions at $s=\frac12$. Ann. of Math. 152 (2000).
\bibitem{soundararajanyoung2013} K. Soundararajan and M. P. Young. The prime geodesic theorem. J. Reine Angew. Math. 676 (2013).
\bibitem{stark1966} H. M. Stark. On the asymptotic density of the $k$ free integers. Proc. Amer. Math. Soc. 17 (1966).
\bibitem{stein1965} C. Stein. Approximation of improper prior measures by prior probability measures. In Bernoulli, Bayes, Laplace Anniversary Volume, Springer (1965).
\bibitem{stepanov1969} S. A. Stepanov. On the number of points of a hyperelliptic curve over a finite prime field. Izv. Akad. Nauk SSSR Ser. Mat. 33 (1969).
\bibitem{steutelthiemann1989} F. W. Steutel and J. G. F. Thiemann. On the independence of integer and fractional parts. Statist. Neerlandica 43 (1989).
\bibitem{stieltjes1894} T. J. Stieltjes. Recherches sur les fractions continues. Ann. Fac. Sci. Toulouse 8 (1894).
\bibitem{stopple2014} J. Stopple. Notes on low discriminants and the generalized Newman conjecture. Funct. Approx. Comment. Math. 51 (2014).
\bibitem{sun2026} W. Sun. Tilting Billingsley's model toward a giant prime: two phase transitions. Preprint, arXiv 2609.11735 (2026).
\bibitem{sunada1986} T. Sunada. L functions in geometry and some applications. In Curvature and Topology of Riemannian Manifolds, Lecture Notes in Math. 1201, Springer (1986).
\bibitem{sutherland2019} A. Sutherland. Sato Tate distributions. Contemp. Math. 740 (2019).
\bibitem{swinnertondyer1973} H. P. F. Swinnerton Dyer. On $\ell$ adic representations and congruences for coefficients of modular forms. In Modular Functions of One Variable III, Lecture Notes in Math. 350, Springer (1973).
\bibitem{sxz2024} X. Sun, S. Xu and Z. Zhuang. Annulus crossing formulae for critical planar percolation. Preprint, arXiv 2410.04767 (2024).
\bibitem{takacs1967} L. Tak\'acs. Combinatorial Methods in the Theory of Stochastic Processes. Wiley (1967).
\bibitem{talbot1979} A. Talbot. The accurate numerical inversion of Laplace transforms. J. Inst. Math. Appl. 23 (1979).
\bibitem{taniguchithorne2013} T. Taniguchi and F. Thorne. Secondary terms in counting functions for cubic fields. Duke Math. J. 162 (2013).
\bibitem{tate1950} J. Tate. Fourier analysis in number fields and Hecke's zeta functions. Thesis, Princeton (1950); in Algebraic Number Theory, Cassels and Fr\"ohlich eds., Academic Press (1967).
\bibitem{tate1965} J. Tate. Algebraic cycles and poles of zeta functions. In Arithmetical Algebraic Geometry, Harper and Row (1965).
\bibitem{tate1967} J. Tate. Fourier analysis in number fields and Hecke's zeta functions. In Algebraic Number Theory, J. W. S. Cassels and A. Fr\"ohlich (eds.), Academic Press (1967).
\bibitem{taylor2008} R. Taylor. Automorphy for some $\ell$ adic lifts of automorphic mod $\ell$ Galois representations II. Publ. Math. IHES 108 (2008).
\bibitem{taylorwiles1995} R. Taylor and A. Wiles. Ring theoretic properties of certain Hecke algebras. Ann. of Math. 141 (1995).
\bibitem{tenenbaum2000} G. Tenenbaum. A rate estimate in Billingsley's theorem for the size distribution of large prime factors. Quart. J. Math. 51 (2000).
\bibitem{tenenbaum2015} G. Tenenbaum. Introduction to Analytic and Probabilistic Number Theory, third edition. Graduate Studies in Mathematics 163, American Mathematical Society (2015).
\bibitem{tenenbaum2017} G. Tenenbaum. Moyennes effectives de fonctions multiplicatives complexes. Ramanujan J. 44 (2017).
\bibitem{terras2011} A. Terras. Zeta Functions of Graphs: A Stroll through the Garden. Cambridge University Press (2011).
\bibitem{thorin1977} O. Thorin. On the infinite divisibility of the Pareto distribution. Scand. Actuar. J. (1977).
\bibitem{thorin1977b} O. Thorin. On the infinite divisibility of the lognormal distribution. Scand. Actuar. J. (1977).
\bibitem{thorin1978} O. Thorin. An extension of the notion of a generalized $\Gamma$ convolution. Scand. Actuar. J. (1978).
\bibitem{thorner2021} J. Thorner. Effective forms of the Sato Tate conjecture. Res. Math. Sci. 8 (2021).
\bibitem{thornerzaman2018} J. Thorner and A. Zaman. A Chebotarev variant of the Brun Titchmarsh theorem and bounds for the Lang Trotter conjectures. Int. Math. Res. Not. IMRN (2018).
\bibitem{titchmarsh1986} E. C. Titchmarsh. The Theory of the Riemann Zeta Function. Second edition, revised by D. R. Heath Brown, Oxford University Press (1986).
\bibitem{tsang1984} K. M. Tsang. The distribution of the values of the Riemann zeta function. PhD thesis, Princeton University (1984).
\bibitem{tsfasmanvladut2002} M. A. Tsfasman and S. G. Vl\u{a}du\c{t}. Infinite global fields and the generalized Brauer Siegel theorem. Mosc. Math. J. 2 (2002).
\bibitem{tunnell1981} J. Tunnell. Artin's conjecture for representations of octahedral type. Bull. Amer. Math. Soc. 5 (1981).
\bibitem{turing1953} A. M. Turing. Some calculations of the Riemann zeta function. Proc. London Math. Soc. 3 (1953).
\bibitem{vandantzig1958} D. van Dantzig. Prize questions. Nieuw Arch. Wisk. 6 (1958).
\bibitem{vonneumann1934} J. von Neumann. Zum Haarschen Mass in topologischen Gruppen. Compositio Math. 1 (1934).
\bibitem{vonneumann1936} J. von Neumann. The uniqueness of Haar's measure. Mat. Sb. (N.S.) 1 (1936).
\bibitem{voronin1975} S. M. Voronin. Theorem on the universality of the Riemann zeta function. Izv. Akad. Nauk SSSR Ser. Mat. 39 (1975).
\bibitem{wald1947} A. Wald. Sequential Analysis. Wiley (1947).
\bibitem{wald1949} A. Wald. Note on the consistency of the maximum likelihood estimate. Ann. Math. Statist. 20 (1949).
\bibitem{waldspurger1981} J. L. Waldspurger. Sur les coefficients de Fourier des formes modulaires de poids demi entier. J. Math. Pures Appl. 60 (1981).
\bibitem{waldspurger1985} J. L. Waldspurger. Sur les valeurs de certaines fonctions L automorphes en leur centre de sym\'etrie. Compositio Math. 54 (1985).
\bibitem{waldwolfowitz1948} A. Wald and J. Wolfowitz. Optimum character of the sequential probability ratio test. Ann. Math. Statist. 19 (1948).
\bibitem{watkins2004} M. Watkins. Real zeros of real odd Dirichlet $L$ functions. Math. Comp. 73 (2004).
\bibitem{webb2015} C. Webb. The characteristic polynomial of a random unitary matrix and Gaussian multiplicative chaos: the $L^2$ phase. Electron. J. Probab. 20 (2015).
\bibitem{weil1940} A. Weil. L'int\'egration dans les groupes topologiques et ses applications. Hermann, Paris (1940).
\bibitem{weil1948} A. Weil. Sur les courbes alg\'ebriques et les vari\'et\'es qui s'en d\'eduisent. Hermann, Paris (1948).
\bibitem{weil1948b} A. Weil. On some exponential sums. Proc. Nat. Acad. Sci. USA 34 (1948).
\bibitem{weil1952} A. Weil. Sur les formules explicites de la th\'eorie des nombres premiers. Comm. S\'em. Math. Univ. Lund, volume d\'edi\'e \`a Marcel Riesz (1952).
\bibitem{weil1961} A. Weil. Adeles and Algebraic Groups. Lecture notes, Institute for Advanced Study, Princeton (1961).
\bibitem{weil1967} A. Weil. \"Uber die Bestimmung Dirichletscher Reihen durch Funktionalgleichungen. Math. Ann. 168 (1967).
\bibitem{weil1972} A. Weil. Sur les formules explicites de la th\'eorie des nombres. Izv. Akad. Nauk SSSR Ser. Mat. 36 (1972).
\bibitem{weisstein} E. W. Weisstein. Langlands Program. MathWorld, A Wolfram Web Resource. https://mathworld.wolfram.com/LanglandsProgram.html.
\bibitem{welchpeers1963} B. L. Welch and H. W. Peers. On formulae for confidence points based on integrals of weighted likelihoods. J. R. Statist. Soc. B 25 (1963).
\bibitem{wesolowski2022} B. Wesolowski. The supersingular isogeny path and endomorphism ring problems are equivalent. In 62nd Annual IEEE Symposium on Foundations of Computer Science, FOCS 2021, IEEE (2022).
\bibitem{weyl1912} H. Weyl. Das asymptotische Verteilungsgesetz der Eigenwerte linearer partieller Differentialgleichungen. Math. Ann. 71 (1912).
\bibitem{weyl1916} H. Weyl. \"Uber die Gleichverteilung von Zahlen mod. Eins. Math. Ann. 77 (1916).
\bibitem{weyl1925} H. Weyl. Theorie der Darstellung kontinuierlicher halbeinfacher Gruppen durch lineare Transformationen I, II, III. Math. Z. 23 (1925), 24 (1926).
\bibitem{widder1941} D. V. Widder. The Laplace Transform. Princeton University Press (1941).
\bibitem{wiener1932} N. Wiener. Tauberian theorems. Ann. of Math. 33 (1932).
\bibitem{wienerikehara1931} S. Ikehara. An extension of Landau's theorem in the analytical theory of numbers. J. Math. Phys. 10 (1931).
\bibitem{wiles1995} A. Wiles. Modular elliptic curves and Fermat's Last Theorem. Ann. of Math. 141 (1995).
\bibitem{williams1974} K. S. Williams. Mertens' theorem for arithmetic progressions. J. Number Theory 6 (1974).
\bibitem{williams1990} D. Williams. Brownian motion and the Riemann zeta function. In Disorder in Physical Systems, Oxford University Press (1990).
\bibitem{wilson1996} D. B. Wilson. Generating random spanning trees more quickly than the cover time. In Proc. 28th ACM Symp. Theory of Computing (1996).
\bibitem{wilton1930} J. R. Wilton. Congruence properties of Ramanujan's function $\tau(n)$. Proc. London Math. Soc. 31 (1930).
\bibitem{wintner1944} A. Wintner. Random factorizations and Riemann's hypothesis. Duke Math. J. 11 (1944).
\bibitem{wirsing1961} E. Wirsing. Das asymptotische Verhalten von Summen \"uber multiplikative Funktionen. Math. Ann. 143 (1961).
\bibitem{wood2017} M. M. Wood. The distribution of sandpile groups of random graphs. J. Amer. Math. Soc. 30 (2017).
\bibitem{wood2019} M. M. Wood. Random integral matrices and the Cohen Lenstra heuristics. Amer. J. Math. 141 (2019).
\bibitem{yamazato1978} M. Yamazato. Unimodality of infinitely divisible distribution functions of class L. Ann. Probab. 6 (1978).
\bibitem{young2006} M. Young. Low lying zeros of families of elliptic curves. J. Amer. Math. Soc. 19 (2006).
\bibitem{zagier1977} D. Zagier. Modular forms whose Fourier coefficients involve zeta functions of quadratic fields. In Modular Functions of One Variable VI, Lecture Notes in Math. 627, Springer (1977).
\bibitem{zarhin2000} Yu. G. Zarhin. Hyperelliptic Jacobians without complex multiplication. Math. Res. Lett. 7 (2000).
\bibitem{zhukleinlukovits1996} H. Y. Zhu, D. J. Klein and I. Lukovits. Extensions of the Wiener number. J. Chem. Inf. Comput. Sci. 36 (1996).
\bibitem{zolotarev1872} G. Zolotarev. Nouvelle d\'emonstration de la loi de r\'eciprocit\'e de Legendre. Nouv. Ann. Math. (2) 11 (1872).
\bibitem{zywina2011} D. Zywina. A refinement of Koblitz's conjecture. Int. J. Number Theory 7 (2011).
\bibitem{zywina2015} D. Zywina. Bounds for the Lang Trotter conjectures. In SCHOLAR: a Scientific Celebration Highlighting Open Lines of Arithmetic Research, Contemp. Math. 655 (2015).
\bibitem{zywina2022} D. Zywina. Explicit open images for elliptic curves over $\mathbb{Q}$. Preprint, arXiv 2206.14959 (2022).
\end{thebibliography}
